% Manuel de l'ingénieur — Simetriik Forge (KOMA-Script, XeLaTeX)
\documentclass[a4paper,11pt,twoside,openany,DIV=calc,parskip=half-,
  headings=big,numbers=noenddot,toc=listof]{scrbook}
\usepackage{fontspec}
\setmainfont{Libertinus Serif}
\setsansfont{Fira Sans}
\setmonofont{Fira Mono}[Scale=MatchLowercase]
\usepackage[french]{babel}
\usepackage{microtype}
\usepackage{xcolor}
\usepackage{graphicx,booktabs,tabularx,colortbl,enumitem,float}
\usepackage[obeyspaces]{url}
\newcolumntype{X}{>{\raggedright\arraybackslash}X}
\usepackage{fancyvrb,fvextra}
\usepackage[skins,breakable]{tcolorbox}
\usepackage{tikz,pifont}
\usetikzlibrary{shapes.geometric,arrows.meta,positioning,calc,fit,backgrounds,decorations.pathmorphing}
\usepackage[hidelinks,pdfusetitle]{hyperref}
\usepackage{scrlayer-scrpage}

\emergencystretch=3em
\raggedbottom

% ---------- Palette ----------
\definecolor{accent}{HTML}{1F4E79}    % bleu enclume (logo)
\definecolor{accentl}{HTML}{5B9BD5}   % bleu clair (logo)
\definecolor{ochre}{HTML}{B7791F}     % fil rouge Sika
\definecolor{leaf}{HTML}{2E7D5B}      % exercices
\definecolor{brick}{HTML}{B03A2E}     % pièges
\definecolor{ink}{HTML}{1C2430}
\definecolor{codebg}{HTML}{F3F5F8}
\definecolor{grey}{HTML}{6B7785}
\color{ink}

% ---------- Typographie KOMA ----------
\setkomafont{disposition}{\sffamily\bfseries\color{accent}}
\setkomafont{chapter}{\sffamily\bfseries\huge\color{accent}}
\setkomafont{section}{\sffamily\bfseries\Large\color{accent}}
\setkomafont{subsection}{\sffamily\bfseries\large\color{accent}}
\setkomafont{descriptionlabel}{\sffamily\bfseries\color{accent}}
\setkomafont{caption}{\sffamily\small}
\setkomafont{captionlabel}{\sffamily\bfseries\small\color{accent}}
\setkomafont{part}{\sffamily\bfseries\Huge\color{accent}}
\setkomafont{partnumber}{\sffamily\bfseries\fontsize{70}{70}\selectfont\color{accent!35}}
\renewcommand*{\chapterformat}{\raisebox{-.15em}{\fontsize{46}{46}\selectfont\color{accentl!70}\thechapter}\enskip}
\renewcommand*{\chapterheadendvskip}{\vspace*{1.2\baselineskip}}
\renewcommand*{\partformat}{\partname~\thepart\autodot}
\setkomafont{pageheadfoot}{\sffamily\small\color{grey}}
\setkomafont{pagenumber}{\sffamily\small\color{accent}}
\KOMAoptions{headsepline=true}
\pagestyle{scrheadings}
\automark[section]{chapter}
\clearpairofpagestyles
\lehead{\pagemark\quad\headmark}
\rohead{\headmark\quad\pagemark}
\lofoot{Simetriik Forge}
\rofoot{Manuel de l'ingénieur}
\setlength{\parindent}{0pt}
\setlist{itemsep=2pt,topsep=3pt}
\renewcommand{\arraystretch}{1.25}
\setcounter{secnumdepth}{2}
\setcounter{tocdepth}{1}

% ---------- Code et commandes ----------
\DeclareUrlCommand\cmd{\urlstyle{tt}\def\UrlFont{\ttfamily\small\color{accent!85!black}}}
\Urlmuskip=0mu plus 1mu\relax
\def\UrlBreaks{\do\/\do\-\do\_\do\.\do\:}
\fvset{fontsize=\small,breaklines,breakanywhere,breaksymbolleft={},frame=none}
\newtcolorbox{code}[1][]{enhanced,breakable,boxrule=0pt,frame hidden,sharp corners,
  colback=codebg,borderline west={2.5pt}{0pt}{accentl},left=3mm,right=2mm,top=1mm,bottom=1mm,
  fontupper=\small,before skip=6pt,after skip=8pt,#1}
\newcommand{\legende}[1]{\par\vspace{-3pt}{\sffamily\scriptsize\color{grey}#1}\par\vspace{4pt}}

% ---------- Pastilles andragogiques ----------
\newcommand{\encadre}[3]{% #1 nom, #2 couleur, #3 libellé
  \newtcolorbox{#1}[1][]{enhanced,colback=#2!5,colframe=#2,boxrule=0pt,
    sharp corners=south,arc=2pt,borderline west={3pt}{0pt}{#2},
    left=3mm,right=3mm,top=2mm,bottom=2mm,before skip=8pt,after skip=10pt,
    title={##1},attach boxed title to top left={xshift=-1mm,yshift=-2mm},
    boxed title style={colback=#2,colframe=#2,sharp corners,left=2mm,right=2mm,top=.5mm,bottom=.5mm},
    fonttitle=\sffamily\bfseries\small\color{white},coltitle=white}}
\encadre{objectifs}{accent}{}
\encadre{situation}{ochre}{}
\encadre{concept}{accentl!80!black}{}
\encadre{piege}{brick}{}
\encadre{vousjouez}{leaf}{}
\encadre{retenir}{accent}{}

% Pictos simples (pas de dépendance à une police d'icônes)
\newcommand{\sika}{\textbf{Sika}}
\newcommand{\q}[1]{\og #1\fg{}}
\newcommand{\FCFA}[1]{#1\,FCFA}
\newcommand{\skill}[1]{\texttt{\color{accent}#1}}
\newcommand{\source}[1]{\par\vspace{1ex}{\raggedright\sffamily\scriptsize\color{grey}Maintenance : à garder synchronisé avec {\ttfamily\detokenize{#1}}.\par}}
\newenvironment{acquis}{\begin{vousjouez}[Vérifiez vos acquis]\begin{enumerate}[label=\sffamily\bfseries\color{leaf}Q\arabic*.,leftmargin=2.2em]}{\end{enumerate}{\small\itshape Réponses en annexe~\ref{app:corriges}.}\end{vousjouez}}

% ---------- Page de titre ----------
\title{Simetriik Forge}
\subtitle{Le manuel de l'ingénieur}
\author{Simetriik Solutions}
\date{Version du plugin : 0.2.1 \textbullet{} \today}
\hypersetup{pdftitle={Simetriik Forge — Manuel de l'ingénieur},pdfauthor={Simetriik Solutions}}

\begin{document}
\frontmatter
\begin{titlepage}
\thispagestyle{empty}
\begin{tikzpicture}[remember picture,overlay]
  \fill[accent] (current page.north west) rectangle ([yshift=-15.5cm]current page.north east);
  \fill[accentl] ([yshift=-15.5cm]current page.north west) rectangle ([yshift=-15.9cm]current page.north east);
  \foreach \i in {0,...,6}{\fill[white,opacity=.07] ([xshift=\i*3cm-1cm,yshift=-15.5cm]current page.north west) -- ++(1.5cm,0) -- ++(-1.5cm,8cm) -- ++(-1.5cm,0) -- cycle;}
\end{tikzpicture}
\vspace*{1.2cm}
\begin{center}
  \includegraphics[width=3.4cm]{../assets/logo-dark.png}\\[1.6cm]
  {\sffamily\bfseries\fontsize{44}{48}\selectfont\color{white}Simetriik Forge}\\[.6cm]
  {\sffamily\LARGE\color{white}Le manuel de l'ingénieur}\\[.5cm]
  {\sffamily\large\color{white!85!accentl}Pratiquer la démarche Lumière, du besoin à la mise en production}
\end{center}
\vfill
{\sffamily\large\color{accent}\textbf{Un fil rouge :} la plateforme \sika{}, qui dématérialise\\
les tontines d'une association de commerçantes, étape après étape.}
\vspace{1.5cm}

{\sffamily\small\color{grey}Simetriik Solutions \textbullet{} Version du plugin 0.2.1 \textbullet{} \today}
\end{titlepage}

\addchap{Comment lire ce manuel}

Ce manuel s'adresse à vous, ingénieur ou ingénieure, qui avez déjà livré des logiciels et
qui découvrez la démarche \emph{Lumière} et son plugin \textbf{Simetriik Forge}. Vous n'avez pas besoin
de tout lire dans l'ordre, ni de tout retenir : il est construit pour que vous pratiquiez vite, sur
un vrai projet, et que vous reveniez aux pages utiles quand le besoin se présente.

\section*{Ce que vous saurez faire}
\begin{itemize}
  \item cadrer un besoin en objectifs mesurables, et refuser une fonctionnalité qui n'en sert aucun ;
  \item bâtir un langage commun entre métier et code, et le faire respecter automatiquement ;
  \item écrire une spécification que le client peut valider et qu'un test peut vérifier ;
  \item consigner une décision d'architecture, chiffrée en FCFA, pour que personne ne la rediscute dans six mois ;
  \item déclencher la recette et la production sans improviser ;
  \item savoir \emph{pourquoi} un garde-fou vous refuse un commit, et quoi faire ensuite.
\end{itemize}

\section*{Un fil rouge : la plateforme Sika}
Tout au long du manuel, un même projet fictif sert d'exemple : \sika{}, une plateforme qui permet aux
groupes de tontine d'une association de commerçantes de cotiser par Mobile Money au lieu de
passer par un cahier et des espèces. Vous le suivrez de la première conversation avec la
présidente jusqu'à la revue post-production. Chaque concept est d'abord défini, puis
illustré sur \sika{}, puis proposé à votre propre projet.

\begin{situation}[Fiction]
\sika{}, l'association, les montants et les noms qui l'entourent sont inventés pour les besoins du
manuel. Les chiffres (coûts, taux, délais) illustrent un raisonnement ; ils ne sont pas des
recommandations et ne doivent pas être recopiés tels quels.
\end{situation}

\section*{Comment un chapitre est construit}
Les adultes apprennent mieux quand ils savent \emph{pourquoi} ils apprennent, partent de leur
expérience et appliquent tout de suite. Chaque chapitre d'étape suit donc le même rythme,
repérable à ses encadrés.

\begin{objectifs}[Objectifs]
Ce que vous serez capable de faire à la fin du chapitre, énoncé en verbes d'action.
\end{objectifs}
\begin{situation}[Sur le terrain : \sika{}]
Une scène du projet fil rouge, pour poser le problème avant la théorie.
\end{situation}
\begin{concept}[Concept]
La définition précise d'un mot de la démarche. Un seul sens, relu à chaque chapitre.
\end{concept}
\begin{piege}[Piège]
Une erreur fréquente, et comment l'éviter.
\end{piege}
\begin{vousjouez}[À vous de jouer]
Un exercice court à refaire sur \emph{votre} projet, tout de suite.
\end{vousjouez}
\begin{retenir}[À retenir]
L'essentiel en quelques lignes, à relire quand vous revenez au chapitre.
\end{retenir}

Chaque chapitre d'étape se termine par quelques questions d'auto-évaluation
(« Vérifiez vos acquis »), dont les réponses commentées figurent en annexe~\ref{app:corriges}.

\section*{Trois parcours de lecture}
\noindent\begin{tabularx}{\linewidth}{>{\raggedright\arraybackslash\bfseries\color{accent}}p{3.2cm}X}
\toprule
Je découvre & Lisez les chapitres \ref{ch:demarche} et \ref{ch:prise} dans l'ordre, puis les étapes une à une. Appliquez chaque « À vous de jouer » avant de passer au chapitre suivant. \\
Je connais Lumière & Parcourez le chapitre \ref{ch:prise}, puis allez à l'étape qui vous bloque. L'annexe~\ref{app:ref} liste tous les skills et scripts. \\
Je suis pressé(e) & Une petite fonctionnalité ? Chapitre \ref{ch:rapide}. Un commit refusé ? Annexe~\ref{app:depannage}. \\
\bottomrule
\end{tabularx}

\section*{Conventions}
\begin{itemize}
  \item Les commandes, chemins et noms de fichiers sont en \texttt{police à chasse fixe}. Les skills s'invoquent par leur nom, par exemple \skill{forge-cadrage}.
  \item Le code, les noms de fichiers et les routes sont en anglais ; les commentaires et la documentation, en français.
  \item Les sorties de scripts reproduites ici ont été obtenues en exécutant réellement les garde-fous sur un projet d'essai.
\end{itemize}

\begin{retenir}[Votre carnet de bord]
Gardez un projet réel sous la main et un carnet ouvert : notez, à chaque « À vous de jouer »,
ce qui s'applique à votre contexte et ce qui ne s'y applique pas. Une démarche ne devient la
vôtre que lorsque vous l'avez adaptée.
\end{retenir}

\tableofcontents

\mainmatter
\part{Comprendre}
\chapter{Pourquoi une démarche}\label{ch:demarche}

\begin{objectifs}[Objectifs]
\begin{itemize}
  \item expliquer en trois phrases pourquoi un projet logiciel dérape ;
  \item citer les six principes de la démarche et leur garde-fou ;
  \item situer chacune des sept étapes et son skill ;
  \item dire ce que le plugin ne fera jamais à votre place.
\end{itemize}
\end{objectifs}

\begin{situation}[Sur le terrain : \sika{}]
Un mardi, la présidente de l'Association Entraide Dantokpa reçoit une équipe de développeurs.
Elle décrit son problème en une phrase : \q{Chaque semaine, quelqu'un dit qu'il a payé et le cahier
dit le contraire.} Dix minutes plus tard, un développeur propose déjà un écran de connexion, une
messagerie interne et un classement des meilleures cotisantes. Personne n'a demandé combien de
cotisations sont contestées aujourd'hui. Ce chapitre explique comment éviter cette scène.
\end{situation}

\section{Le problème}
Un projet logiciel dérape rarement à cause du code. Il dérape parce qu'on construit la mauvaise
chose, parce qu'on se comprend mal, ou parce qu'on décide trop tôt sans le dire. La démarche
Lumière répond à ces trois risques par une chaîne d'étapes courtes ; chacune produit un
livrable \emph{validé par un humain} avant que la suivante ne commence.

Le plugin \textbf{Simetriik Forge} met cette démarche en pratique avec un assistant de code. Il
apporte deux choses, à ne pas confondre :
\begin{description}
  \item[Des skills] des instructions que l'assistant suit pour une étape (poser les bonnes questions, remplir le bon gabarit). Ils \emph{guident}.
  \item[Des garde-fous] de petits scripts (\cmd{scripts/}) qui \emph{échouent} quand une règle est violée. Ils \emph{contrôlent}, sans interprétation, et fonctionnent aussi en intégration continue.
\end{description}

\begin{concept}[Garde-fou]
Un script déterministe qui vérifie une règle de la démarche et retourne un code d'échec si elle
est violée. Il ne juge pas la qualité : il constate un fait (un statut, un nom, un fichier
présent). La qualité reste l'affaire de la relecture humaine.
\end{concept}

\section{Les six principes}
\noindent\begin{tabularx}{\linewidth}{>{\raggedright\arraybackslash}p{6.3cm}X}
\toprule
\textbf{Principe} & \textbf{Ce qui le fait respecter} \\ \midrule
Pas de code sans spécification validée. La validation est humaine. & \cmd{check-spec-validated.sh}, hook \emph{pre-commit} \\
Chaque fonctionnalité remonte à un objectif mesurable. & l'impact map, colonne « Objectif » obligatoire \\
Un terme = une définition = un nom technique. & \cmd{check-glossary-terms.sh} (FF-05) \\
Le contrat d'API précède le code. & \cmd{check-openapi-first.sh} \\
Les décisions structurantes sont consignées, stack comprise. & \cmd{check-adr-present.sh}, ADR \\
Ce qui s'automatise s'automatise. & les scripts, en local et en CI \\
\bottomrule
\end{tabularx}

\section{Les sept étapes d'un coup d'œil}
La démarche compte sept étapes numérotées de 0 à 6. La mise en œuvre du code, quant à elle,
se glisse entre le contrat et la recette.

\begin{figure}[H]
\centering
\resizebox{\linewidth}{!}{%
\begin{tikzpicture}[
  st/.style={draw=accent,fill=accent!7,rounded corners=3pt,minimum width=3.1cm,minimum height=15mm,align=center,font=\sffamily\small,text=ink},
  co/.style={draw=ochre,fill=ochre!10,rounded corners=3pt,minimum width=3.1cm,minimum height=15mm,align=center,font=\sffamily\small,text=ink},
  num/.style={circle,fill=accent,text=white,font=\sffamily\bfseries\footnotesize,inner sep=2pt},
  ar/.style={-{Stealth[length=2.2mm]},thick,accent}]
  \node[st] (a) at (0,0)   {\textbf{Cadrage}\\[-1pt]{\scriptsize forge-cadrage}};
  \node[st] (b) at (3.8,0) {\textbf{Domaine}\\[-1pt]{\scriptsize forge-domaine}};
  \node[st] (c) at (7.6,0) {\textbf{Spécification}\\[-1pt]{\scriptsize forge-spec}};
  \node[st] (d) at (11.4,0){\textbf{Architecture}\\[-1pt]{\scriptsize forge-architecture}};
  \node[st] (e) at (11.4,-2.6){\textbf{Contrat d'API}\\[-1pt]{\scriptsize forge-contrat}};
  \node[co] (f) at (7.6,-2.6){\textbf{Code}\\[-1pt]{\scriptsize forge-implemente}};
  \node[st] (g) at (3.8,-2.6){\textbf{Recette}\\[-1pt]{\scriptsize forge-recette}};
  \node[st] (h) at (0,-2.6){\textbf{Production}\\[-1pt]{\scriptsize forge-prod}};
  \node[num] at (a.north west) {0}; \node[num] at (b.north west) {1}; \node[num] at (c.north west) {2};
  \node[num] at (d.north west) {3}; \node[num] at (e.north west) {4}; \node[num] at (g.north west) {5}; \node[num] at (h.north west) {6};
  \draw[ar] (a)--(b); \draw[ar] (b)--(c); \draw[ar] (c)--(d); \draw[ar] (d)--(e);
  \draw[ar] (e)--(f); \draw[ar] (f)--(g); \draw[ar] (g)--(h);
  \draw[ar,dashed,grey] (h.south) -- ++(0,-8mm) -| node[pos=.25,below,font=\sffamily\scriptsize,text=grey]{lot suivant : retour au cadrage} ([xshift=-8mm]a.south west) |- (a.west);
\end{tikzpicture}}
\caption{Le chemin de la démarche : sept étapes validées, et le code entre le contrat et la recette.}
\end{figure}

À côté de ce chemin, trois skills transversaux : \skill{forge-init} (démarrer), \skill{forge-statut}
(où en suis-je ?) et \skill{forge-rapide} (le cycle court des petites fonctionnalités, chapitre~\ref{ch:rapide}).

\section{Ce que le plugin ne fait pas}
Il ne déploie jamais en production de lui-même, ne valide aucune spécification, ne pousse aucun
code sans demande et n'invente aucune règle métier : quand une information manque, il pose la
question. C'est volontaire. La validation engage une responsabilité (celle du client, ou la
vôtre) ; elle ne se délègue pas à un outil.

\begin{piege}[Piège : « l'assistant l'a validé »]
Un assistant peut passer un statut à \emph{Validée} en une ligne ; le script, lui, le croira.
C'est pourquoi le plugin lui interdit de le faire : la ligne \emph{Validé par} porte un nom et
une date \emph{humains}. Une spec validée par personne est une spec non validée.
\end{piege}

\begin{vousjouez}[À vous de jouer]
Prenez un projet que vous avez vécu et qui a dérapé. En une phrase chacune, laquelle de ces trois causes
l'explique le mieux : \emph{mauvaise chose construite}, \emph{malentendu}, \emph{décision tardive non dite} ?
Notez-la dans votre carnet : vous y reviendrez aux chapitres~\ref{ch:cadrage}, \ref{ch:domaine} et \ref{ch:archi}.
\end{vousjouez}

\begin{retenir}[À retenir]
\begin{itemize}
  \item Trois risques : mauvaise chose, malentendu, décision tardive.
  \item Les \emph{skills} guident ; les \emph{garde-fous} contrôlent ; l'\emph{humain} valide.
  \item Sept étapes, un livrable validé chacune, le code entre le contrat et la recette.
\end{itemize}
\end{retenir}
\source{README.md}

\begin{acquis}
  \item Quelle différence y a-t-il entre un skill et un garde-fou ?
  \item Qui peut passer une spécification au statut \emph{Validée} ?
  \item Dans quel ordre viennent le contrat d'API et le code ?
\end{acquis}

\chapter{Prise en main}
\section{Installation}
\begin{itemize}
  \item \textbf{Claude Code} : \cmd{claude --plugin-dir /chemin/vers/simetriik-forge}.
  \item \textbf{Codex et OpenCode} : lancer \cmd{./install.sh} (liens symboliques dans \cmd{~/.agents/skills}), puis redémarrer l'outil si les skills n'apparaissent pas.
\end{itemize}
Les skills s'invoquent par leur nom, par exemple \cmd{/simetriik-forge:forge-init} sous Claude Code.

\section{Démarrer un projet : \texttt{forge-init}}
Dans un dossier vide, lancez \cmd{forge-init}. Le skill :
\begin{enumerate}
  \item vérifie qu'il n'existe pas déjà de \cmd{docs/} ni de \cmd{AGENTS.md} (il n'écrase jamais) ;
  \item vous pose quatre questions, une à la fois : nom du projet, institution cliente, chef de projet, date de démarrage ;
  \item copie le dossier de conception (\cmd{docs/}, \cmd{api/}, \cmd{AGENTS.md}, \cmd{.github/}) et les scripts de garde-fou dans \cmd{scripts/} ;
  \item remplace les marqueurs \cmd{[NOM DU PROJET]}, \cmd{[INSTITUTION]}, \cmd{[NOM]}, \cmd{[JJ/MM/AAAA]} ;
  \item crée \cmd{CLAUDE.md} contenant \cmd{@AGENTS.md} ;
  \item annonce la suite : \cmd{forge-cadrage}.
\end{enumerate}
\astuce{La stack technique n'est \emph{pas} demandée à l'initialisation : c'est une décision d'architecture, prise à l'étape~3 par un ADR, à partir du cadrage et du domaine.}

\section{Structure du dossier de conception}
\begin{lstlisting}
docs/00-cadrage/        vision, impact map
docs/01-domaine/        event storming, glossaire, context map
docs/02-specs/          SPEC-NNN.md, features/*.feature, exigences non fonctionnelles
docs/03-architecture/   adr/, c4/workspace.dsl, fitness-functions.md
docs/04-livrables-client/
docs/05-recette/        plan-recette.md
docs/06-production/     deploiement, runbook, rollback, monitoring
api/openapi.yaml        contrat d'API
AGENTS.md               consignes pour les assistants de code
\end{lstlisting}

\section{Suivre sa progression}
Deux outils, tous deux en lecture seule :
\begin{itemize}
  \item \cmd{forge-statut} : tableau de bord complet (état de chaque étape et de chaque spec, garde-fous en échec) et \emph{une seule} prochaine action recommandée.
  \item \cmd{scripts/forge-progress.sh} : barre de progression sur les sept étapes, affichée à la fin de chaque skill, par exemple
\par\smallskip{\small\cmd{Progression [███░░░░] 3/7 · étape en cours : Architecture · suite : forge-contrat}}\par
\end{itemize}
Une étape compte comme faite lorsque son livrable est au statut \emph{Validée} (ADR : \emph{Acceptée}).
Pour l'architecture, la barre exige au moins deux ADR acceptés (l'ADR-0001 du monolithe modulaire et l'ADR de stack).

\section{Les statuts}
Brouillon $\rightarrow$ En revue $\rightarrow$ \textbf{Validée}. Seul un humain passe un livrable à \emph{Validée} (en indiquant son nom et la date).
\source{skills/forge-init, forge-statut}

\part{Pratiquer, étape par étape}
\chapter{Cadrer le besoin}\label{ch:cadrage}

\begin{objectifs}[Objectifs --- étape 0, \texttt{forge-cadrage}]
\begin{itemize}
  \item formuler une vision produit et des objectifs \emph{mesurables} ;
  \item construire une impact map reliant pourquoi, qui, comment et quoi ;
  \item prioriser avec MoSCoW et écarter ce qui ne sert aucun objectif ;
  \item reconnaître un chiffre inconnu et le noter \q{À mesurer} plutôt que de l'inventer.
\end{itemize}
\end{objectifs}

\begin{situation}[Sur le terrain : \sika{}]
Retour chez la présidente. Cette fois, avant tout écran, on pose des questions : combien de
groupes ? combien de membres ? combien de cotisations sont contestées chaque semaine ? La
trésorière générale a compté sur trois mois : 12\,\% des cotisations ont été discutées. Elle ne sait en
revanche pas combien de temps s'écoule entre la dernière cotisation et la remise du pot. Deux
objectifs émergent --- l'un avec un chiffre, l'autre sans.
\end{situation}

\section{Le concept : relier chaque fonctionnalité à un objectif}
\begin{concept}[Objectif mesurable]
Un objectif qui contient un \emph{indicateur}, une \emph{valeur actuelle}, une \emph{cible} et une
\emph{échéance}. \q{Améliorer la confiance} n'est pas mesurable ; \q{passer de 12\,\% à moins de 2\,\% de
cotisations contestées en six mois} l'est.
\end{concept}

\begin{concept}[Impact map]
Une carte à quatre niveaux qui relie le \emph{pourquoi} (l'objectif) au \emph{quoi} (la fonctionnalité)
par le \emph{qui} (l'acteur) et le \emph{comment} (le changement de comportement visé). Sa règle
d'or : une fonctionnalité qui ne remonte à aucun objectif sort du périmètre.
\end{concept}

\section{Pas à pas avec \texttt{forge-cadrage}}
Le skill vous interroge \emph{une question à la fois}, dans cet ordre :
\begin{enumerate}
  \item pour qui, et quel problème ;
  \item la situation actuelle, le cadre réglementaire (APDP, OHADA, marchés publics), les contraintes terrain (connectivité, terminaux, langues) ;
  \item les objectifs mesurables : indicateur, valeur actuelle, cible, échéance ;
  \item les acteurs, puis les changements de comportement ;
  \item enfin les fonctionnalités.
\end{enumerate}
Il renseigne ensuite \cmd{vision.md}, puis l'impact map (un diagramme et un tableau de priorisation).
Le livrable passe \emph{En revue} ; la validation revient au \textbf{sponsor client}, pas à l'ingénieur.

\section{Ce que \sika{} produit}
\subsection*{La vision}
\begin{code}
\begin{Verbatim}
Pour les groupes de tontine de l'Association Entraide
Dantokpa qui rapprochent chaque semaine un cahier et des
espèces, Sika est une plateforme de cotisation par Mobile
Money qui rend chaque cotisation vérifiable. Contrairement
au cahier, chaque cotisation y est horodatée et confirmée
par l'opérateur.
\end{Verbatim}
\end{code}

\subsection*{Les objectifs}
\noindent\begin{tabularx}{\linewidth}{lXlll}
\toprule
\textbf{\#} & \textbf{Objectif et indicateur} & \textbf{Actuel} & \textbf{Cible} & \textbf{Échéance} \\ \midrule
O1 & Rendre les cotisations indiscutables : \% de cotisations contestées & 12\,\% & $<2\,\%$ & 6 mois \\
O2 & Remettre le pot plus vite : délai de remise du pot & \textit{À mesurer} & $\le 24$\,h & 6 mois \\
\bottomrule
\end{tabularx}

\begin{piege}[Piège : inventer le chiffre qui manque]
Pour O2, personne ne connaît la valeur actuelle. Écrire \q{3 jours} \emph{parce que ça semble juste}
fausse toute la mesure future. Le skill note \emph{À mesurer} ; la première tâche du projet sera de
mesurer.
\end{piege}

\subsection*{L'impact map}
\begin{figure}[H]
\centering
\resizebox{\linewidth}{!}{%
\begin{tikzpicture}[
  n/.style={draw,rounded corners=3pt,align=center,font=\sffamily\small,minimum height=11mm,text width=3.2cm,inner sep=3pt},
  g/.style={n,fill=accent,draw=accent,text=white},
  a/.style={n,fill=accentl!20,draw=accentl},
  i/.style={n,fill=ochre!12,draw=ochre},
  d/.style={n,fill=leaf!12,draw=leaf},
  ar/.style={-{Stealth},thick,grey},
  h/.style={font=\sffamily\bfseries\scriptsize,text=grey}]
  \node[h] at (0,2.3) {POURQUOI}; \node[h] at (4.3,2.3) {QUI}; \node[h] at (8.6,2.3) {COMMENT}; \node[h] at (12.9,2.3) {QUOI};
  \node[g] (o1) at (0,0.9) {O1\\cotisations contestées $<2\,\%$};
  \node[g] (o2) at (0,-1.9) {O2\\pot remis en $\le 24$\,h};
  \node[a] (m) at (4.3,1.3) {Membre};
  \node[a] (t) at (4.3,-1.3) {Trésorière de groupe};
  \node[i] (i1) at (8.6,1.3) {Cotise depuis son téléphone, avec preuve};
  \node[i] (i2) at (8.6,-1.3) {Suit les cotisations sans pointer le cahier};
  \node[d] (d1) at (12.9,1.7) {Cotiser par Mobile Money};
  \node[d] (d2) at (12.9,0.4) {Solde du pot en temps réel};
  \node[d] (d3) at (12.9,-1.9) {Remettre le pot au bénéficiaire};
  \draw[ar] (o1)--(m); \draw[ar] (o1)--(t); \draw[ar] (o2)--(t);
  \draw[ar] (m)--(i1); \draw[ar] (t)--(i2);
  \draw[ar] (i1)--(d1); \draw[ar] (i1)--(d2); \draw[ar] (i2)--(d3);
\end{tikzpicture}}
\caption{L'impact map de \sika{} : chaque fonctionnalité remonte à un objectif.}
\end{figure}

\subsection*{La priorisation MoSCoW}
\begin{concept}[MoSCoW]
Quatre niveaux de priorité : \textbf{M}ust (indispensable au lot), \textbf{S}hould (important, pas bloquant),
\textbf{C}ould (utile si le temps le permet), \textbf{W}on't (explicitement hors lot).
\end{concept}
\noindent\begin{tabularx}{\linewidth}{Xcll}
\toprule
\textbf{Fonctionnalité} & \textbf{Objectif} & \textbf{Priorité} & \textbf{Lot} \\ \midrule
Cotiser par Mobile Money & O1 & Must & Lot 1 \\
Solde du pot en temps réel & O1 & Should & Lot 1 \\
Remettre le pot au bénéficiaire & O2 & Must & Lot 2 \\
Rappel SMS avant l'échéance & O1 & Could & Lot 2 \\
Messagerie vocale intégrée & \textit{aucun} & Won't & hors périmètre \\
\bottomrule
\end{tabularx}
La messagerie vocale, séduisante, ne sert ni O1 ni O2 : le skill la signale \emph{hors périmètre}. C'est exactement
ce que cette étape est faite pour détecter.

\begin{vousjouez}[À vous de jouer]
Sur votre projet : écrivez deux objectifs au format indicateur / actuel / cible / échéance. Pour chaque
chiffre que vous ne connaissez pas, écrivez \emph{À mesurer} et nommez qui peut le mesurer. Dessinez
ensuite une impact map de cinq nœuds. Y a-t-il une fonctionnalité \q{évidente} qui n'en fait pas partie ?
\end{vousjouez}

\begin{retenir}[À retenir]
\begin{itemize}
  \item Un objectif = indicateur + valeur actuelle + cible + échéance ; sinon il n'est pas mesurable.
  \item Pas d'objectif \texttt{O\#}, pas de fonctionnalité.
  \item Inconnu $\Rightarrow$ \emph{À mesurer}, jamais un chiffre inventé.
  \item Le sponsor client valide le cadrage.
\end{itemize}
\end{retenir}
\source{skills/forge-cadrage/SKILL.md}

\begin{acquis}
  \item Pourquoi \q{améliorer l'expérience des membres} n'est-il pas un objectif exploitable ?
  \item Que fait le skill d'une fonctionnalité sans objectif \texttt{O\#} ?
  \item Qui valide le cadrage ?
\end{acquis}

\chapter{Comprendre le domaine}\label{ch:domaine}

\begin{objectifs}[Objectifs --- étape 1, \texttt{forge-domaine}]
\begin{itemize}
  \item mener ou restituer un Event Storming et repérer les points chauds ;
  \item écrire un glossaire où chaque terme a une définition et un nom technique uniques ;
  \item classer les sous-domaines et dessiner une context map ;
  \item faire respecter le vocabulaire par un script, pas par la bonne volonté.
\end{itemize}
\end{objectifs}

\begin{situation}[Sur le terrain : \sika{}]
Atelier d'une demi-journée avec la trésorière générale et trois trésorières de groupe. Au bout d'une
heure, un désaccord surgit : pour l'une, \q{le pot} est la somme collectée ; pour l'autre, c'est la
remise elle-même ; une troisième dit \q{la cagnotte}. Si le code nomme \cmd{Pot}, \cmd{Jackpot} et
\cmd{Fund} trois choses identiques, le bogue est déjà écrit.
\end{situation}

\section{Event Storming : voir le métier se dérouler}
\begin{concept}[Event Storming]
Un atelier où le métier et la technique posent sur un mur, dans l'ordre chronologique, les
\emph{événements} qui se produisent (au participe passé : \q{Cotisation confirmée}), puis les
commandes qui les déclenchent, les acteurs, les politiques (\q{dès que\dots, alors\dots}), les systèmes
externes et, en rouge, les \emph{points chauds}, c'est-à-dire les questions sans réponse.
\end{concept}

\noindent\begin{tabularx}{\linewidth}{lXX}
\toprule
\textbf{Couleur} & \textbf{Élément} & \textbf{Exemple Sika} \\ \midrule
Orange & Événement métier & Cotisation confirmée \\
Bleu & Commande & Cotiser pour le tour \\
Jaune & Acteur & Trésorière de groupe \\
Lilas & Politique & Dès que tous ont cotisé, clôturer le tour \\
Rose & Système externe & Plateforme Mobile Money \\
Rouge & Point chaud & Que faire d'un membre qui ne cotise pas ? \\
\bottomrule
\end{tabularx}

\begin{figure}[H]
\centering
\resizebox{\linewidth}{!}{%
\begin{tikzpicture}[
  ev/.style={fill=ochre!35,draw=ochre,rounded corners=2pt,align=center,font=\sffamily\scriptsize,text width=2.4cm,minimum height=9mm},
  cde/.style={fill=accentl!35,draw=accentl,rounded corners=2pt,align=center,font=\sffamily\scriptsize,text width=2.4cm,minimum height=9mm},
  po/.style={fill=violet!15,draw=violet!60,rounded corners=2pt,align=center,font=\sffamily\scriptsize,text width=2.4cm,minimum height=9mm},
  hot/.style={fill=brick!20,draw=brick,rounded corners=2pt,align=center,font=\sffamily\scriptsize,text width=3.6cm},
  ar/.style={-{Stealth},grey,thick}]
  \node[ev] (e0) at (0,0) {Tour ouvert};
  \node[cde] (c1) at (3.2,0) {Cotiser};
  \node[ev] (e1) at (6.4,0) {Cotisation initiée};
  \node[ev] (e2) at (9.6,0.9) {Cotisation confirmée};
  \node[ev] (e3) at (9.6,-0.9) {Cotisation échouée};
  \node[po] (p1) at (12.8,0.9) {Dès que tous ont cotisé : clôturer};
  \node[ev] (e4) at (16,0.9) {Tour clôturé};
  \node[cde] (c2) at (16,-0.9) {Remettre le pot};
  \node[ev] (e5) at (12.8,-0.9) {Pot remis};
  \draw[ar] (e0)--(c1); \draw[ar] (c1)--(e1); \draw[ar] (e1)--(e2); \draw[ar] (e1)--(e3);
  \draw[ar] (e2)--(p1); \draw[ar] (p1)--(e4); \draw[ar] (e4)--(c2); \draw[ar] (c2)--(e5);
  \node[hot] at (9.6,-2.7) {\textbf{H2} Au bout de combien de temps une cotisation sans réponse est-elle \q{échouée} ?};
  \node[hot] at (14.4,-2.7) {\textbf{H1} Que fait-on d'un membre qui ne cotise pas avant la date limite ?};
\end{tikzpicture}}
\caption{Le flux principal de \sika{}, tel qu'il sort de l'atelier. Les deux points chauds sont en rouge.}
\end{figure}

\begin{piege}[Piège : trancher un point chaud sans son responsable]
H1 touche au règlement intérieur ; seule la présidente peut y répondre. Le skill ne tranche \emph{jamais}
un point chaud lui-même : il l'inscrit dans la table avec un responsable et une échéance. H2, lui,
conditionnera la règle R2 de la spécification (chapitre~\ref{ch:spec}).
\end{piege}

\section{Le glossaire : un seul mot par chose}
\begin{concept}[Langage omniprésent]
Un vocabulaire unique, partagé par le métier et le code. \emph{Un terme} = \emph{une définition} =
\emph{un nom technique anglais}. Les synonymes courants sont listés pour être \textbf{bannis}.
\end{concept}

Voici le glossaire de \sika{} (extrait). La dernière colonne alimente un script.
\noindent\begin{tabularx}{\linewidth}{llX>{\raggedright\arraybackslash}p{3.4cm}}
\toprule
\textbf{Terme} & \textbf{Nom technique} & \textbf{Définition} & \textbf{Synonymes bannis} \\ \midrule
Membre & \cmd{Member} & Personne qui cotise dans un groupe & adhérente, user, client \\
Groupe & \cmd{SavingsGroup} & Ensemble de membres qui cotisent ensemble & cercle \\
Tour & \cmd{Round} & Période au bout de laquelle un seul membre reçoit le pot & cycle, session \\
Cotisation & \cmd{Contribution} & Somme versée par un membre pour un tour & versement, payment \\
Pot & \cmd{Pot} & Somme des cotisations confirmées d'un tour & cagnotte, caisse \\
Bénéficiaire & \cmd{Beneficiary} & Membre qui reçoit le pot d'un tour & gagnant, winner \\
\bottomrule
\end{tabularx}

\subsection*{Le garde-fou FF-05 en action}
Un développeur écrit, par habitude, une classe \cmd{Payment}. Le script lit la colonne des synonymes et cherche
chaque mot dans le code (sortie réelle) :
\begin{code}
\begin{Verbatim}
$ scripts/check-glossary-terms.sh docs/01-domaine/glossaire.md src/
src/a.ts:1:class Payment { amountXof: number }
KO: terme banni « payment »
\end{Verbatim}
\end{code}
Après correction en \cmd{Contribution} :
\begin{code}
\begin{Verbatim}
OK: aucun terme banni
\end{Verbatim}
\end{code}

\begin{piege}[Piège : affaiblir le glossaire pour faire passer le contrôle]
Quand le script crie, la mauvaise réponse est de retirer le synonyme de la liste. La bonne :
renommer dans le code. Le glossaire est un contrat avec le métier.
\end{piege}

\section{Sous-domaines et context map}
\begin{concept}[Bounded context]
Une frontière à l'intérieur de laquelle un modèle et son vocabulaire sont cohérents. Hors de la
frontière, le même mot peut signifier autre chose.
\end{concept}
Les sous-domaines se classent en trois familles : le \textbf{cœur} (ce qui différencie, là où va l'effort de
conception), le \textbf{support} (nécessaire, pas différenciant) et le \textbf{générique} (commun à tous les systèmes, on
l'achète ou on le fait au plus simple).

\begin{figure}[H]
\centering
\resizebox{.95\linewidth}{!}{%
\begin{tikzpicture}[
  bc/.style={draw,rounded corners=4pt,align=center,font=\sffamily\small,minimum width=3.4cm,minimum height=12mm},
  core/.style={bc,fill=accent,draw=accent,text=white},
  sup/.style={bc,fill=ochre!15,draw=ochre},
  gen/.style={bc,fill=grey!15,draw=grey},
  rel/.style={-{Stealth},thick,accent},
  lab/.style={font=\sffamily\scriptsize,text=accent,fill=white,inner sep=1.5pt}]
  \node[core] (c) at (0,0) {\textbf{Cotisation}\\{\scriptsize cœur : tours, cotisations, pot}};
  \node[sup] (n) at (5.6,1.8) {\textbf{Notifications}\\{\scriptsize support : SMS}};
  \node[gen] (i) at (-5.6,1.8) {\textbf{Identité}\\{\scriptsize générique}};
  \node[gen] (p) at (0,-3) {\textbf{Paiement mobile}\\{\scriptsize générique : agrégateur Mobile Money}};
  \draw[rel] (c)--(n) node[midway,lab]{Customer / Supplier};
  \draw[rel] (i)--(c) node[midway,lab]{Conformist};
  \draw[rel] (p)--(c) node[midway,lab,xshift=0pt]{ACL};
\end{tikzpicture}}
\caption{La context map de \sika{} : un cœur, un support, deux contextes génériques.}
\end{figure}

\noindent\begin{tabularx}{\linewidth}{>{\raggedright\arraybackslash}p{3.4cm}>{\raggedright\arraybackslash}p{3.8cm}X}
\toprule
\textbf{Pattern} & \textbf{Entre} & \textbf{Pourquoi} \\ \midrule
ACL (couche anticorruption) & Paiement mobile $\to$ Cotisation & isoler les particularités de chaque opérateur ; le cœur ne parle jamais MTN ou Moov directement \\
Conformist & Identité $\to$ Cotisation & on adopte le modèle de la solution d'identité sans le déformer \\
Customer / Supplier & Cotisation $\to$ Notifications & le cœur exprime ses besoins ; le support s'y conforme \\
\bottomrule
\end{tabularx}

\begin{vousjouez}[À vous de jouer]
Listez dix événements de votre projet, en quelques minutes, au participe passé. Repérez un mot qu'un collègue
et vous employez pour deux choses différentes (ou deux mots pour une chose). Écrivez la ligne de glossaire correspondante, avec
son nom technique anglais et deux synonymes à bannir.
\end{vousjouez}

\begin{retenir}[À retenir]
\begin{itemize}
  \item Un point chaud n'est jamais tranché sans son responsable.
  \item Un terme = une définition = un nom technique ; les synonymes sont bannis et contrôlés (FF-05).
  \item L'effort de conception va au cœur ; le reste s'achète ou reste simple.
\end{itemize}
\end{retenir}
\source{skills/forge-domaine/SKILL.md, scripts/check-glossary-terms.sh}

\begin{acquis}
  \item Pourquoi un événement s'écrit-il au participe passé ?
  \item Que contient la colonne « Synonymes à éviter », et qui la lit ?
  \item Dans \sika{}, pourquoi le paiement mobile est-il isolé derrière une ACL ?
\end{acquis}

\chapter{Spécifier une fonctionnalité}\label{ch:spec}

\begin{objectifs}[Objectifs --- étape 2, \texttt{forge-spec}]
\begin{itemize}
  \item rédiger une spécification complète, règle par règle, avec la source de chaque règle ;
  \item écrire des scénarios Gherkin : nominal, limite, erreur ;
  \item distinguer ce que le script contrôle (le statut) de ce que vous devez relire ;
  \item obtenir une validation humaine, nommée et datée.
\end{itemize}
\end{objectifs}

\begin{situation}[Sur le terrain : \sika{}]
Première fonctionnalité du lot 1 : \q{Cotiser par Mobile Money}. Le développeur est tenté de commencer par
l'écran. La spec va d'abord forcer trois questions : quel montant est accepté ? quand une cotisation
compte-t-elle vraiment ? que se passe-t-il si un membre cotise deux fois ?
\end{situation}

\section{Le concept : une règle, une source, un exemple}
\begin{concept}[Spécification]
Le document qui décrit \emph{ce que la fonctionnalité doit faire}, sans dire comment. Chaque règle métier
y figure avec sa \textbf{source} (texte réglementaire, règlement intérieur, atelier) : une règle sans
source est une opinion.
\end{concept}
\begin{concept}[Gherkin]
Un langage d'exemples en français (\emph{Étant donné, Quand, Alors}). Un fichier \cmd{.feature}
contient une \cmd{Règle} par règle métier, et pour chacune au moins un cas nominal, un cas limite et un cas d'erreur.
Chaque scénario sert à la fois de critère d'acceptation et de test.
\end{concept}

\section{Pas à pas avec \texttt{forge-spec}}
\begin{enumerate}
  \item Le skill lit le glossaire, l'impact map et le gabarit, puis numérote la spec : \cmd{SPEC-001}.
  \item Il copie le gabarit en \cmd{docs/02-specs/SPEC-001.md}, statut \emph{Brouillon}.
  \item Il remplit section par section, une question à la fois : problème, récit, règles métier (chacune avec sa source), données (les données personnelles sont marquées), interfaces, hors périmètre.
  \item Il écrit \cmd{features/SPEC-001.feature}.
  \item Il passe la spec \emph{En revue} ; la validation, métier puis technique, est humaine.
\end{enumerate}
Le skill n'invente jamais une règle : quand elle manque, il l'inscrit en \emph{Questions ouvertes} et vous la demande.

\section{La spécification de \sika{} (extrait)}
L'extrait est condensé pour la lecture. Dans le fichier réel, l'en-tête est un tableau Markdown (\cmd{| Statut | Brouillon |}) : c'est cette ligne, et elle seule, que lit le garde-fou.

\begin{code}
\begin{Verbatim}
# SPEC-001 — Cotiser par Mobile Money

Statut : En revue · Contexte : Cotisation · Objectif : O1
Priorité : Must

## 2. Récit utilisateur
En tant que membre d'un groupe, je veux cotiser depuis mon
téléphone afin que ma cotisation soit enregistrée sans
contestation.

## 3. Règles métier
R1  Le montant est celui fixé pour le groupe.
    Source : règlement intérieur, art. 7
R2  Une cotisation ne compte qu'une fois confirmée par
    l'opérateur.
    Source : atelier Event Storming, point chaud H2
R3  Un membre cotise une seule fois par tour.
    Source : règlement intérieur, art. 8

## 5. Données
amountXof  entier, obligatoire, non personnelle
phone      texte, obligatoire, DONNÉE PERSONNELLE (APDP)

## 9. Questions ouvertes
Q1  Après combien de minutes sans réponse de l'opérateur une
    cotisation est-elle déclarée échouée ?
    Responsable : trésorière générale
\end{Verbatim}
\end{code}

\begin{piege}[Piège : valider avec une question ouverte bloquante]
Tant que Q1 n'a pas de réponse, R2 est incomplète : le cas limite \q{l'opérateur ne répond pas} ne peut
pas s'écrire. Le script, lui, ne le verra pas : il ne lit que le statut. C'est à vous de relire avant
de proposer \emph{Validée}. Ici, la trésorière générale répond \q{15 minutes} ; la question est close, et
le scénario limite peut être écrit.
\end{piege}

\subsection*{Les scénarios Gherkin}
\begin{code}
\begin{Verbatim}
# language: fr
Fonctionnalité: Cotiser par Mobile Money
  En tant que membre d'un groupe
  Je veux cotiser depuis mon téléphone
  Afin que ma cotisation soit enregistrée sans contestation

  Règle: R1 — le montant est celui fixé pour le groupe
    Scénario: Cotisation au montant du groupe
      Étant donné un groupe dont la cotisation est de 5000 FCFA
      Et un membre de ce groupe avec un tour ouvert
      Quand le membre cotise 5000 FCFA
      Alors la cotisation est initiée auprès de l'opérateur
    Scénario: Montant différent de celui du groupe
      Quand le membre cotise 3000 FCFA
      Alors l'action est refusée avec le motif "Montant incorrect"

  Règle: R2 — une cotisation ne compte qu'une fois confirmée
    Scénario: Confirmation par l'opérateur
      Étant donné une cotisation initiée de 5000 FCFA
      Quand l'opérateur confirme la cotisation
      Alors le solde du pot augmente de 5000 FCFA
      Et le membre reçoit un SMS de confirmation
    Scénario: Opérateur sans réponse après 15 minutes
      Étant donné une cotisation initiée depuis 15 minutes
      Quand le délai expire
      Alors la cotisation est échouée et ne compte pas dans le pot

  Règle: R3 — une seule cotisation par membre et par tour
    Scénario: Seconde cotisation dans le même tour
      Étant donné un membre ayant déjà cotisé pour ce tour
      Quand le membre cotise à nouveau
      Alors l'action est refusée avec le motif "Déjà cotisé pour ce tour"
\end{Verbatim}
\end{code}
Remarquez la lecture : le cas limite (15 minutes), le cas nominal et le cas d'erreur vivent côte à côte, sous la règle qu'ils éprouvent.
Remarquez aussi le vocabulaire : \emph{cotisation}, \emph{pot}, \emph{tour} --- pas \q{paiement}, \q{cagnotte}, \q{cycle}.

\section{Le garde-fou : le statut, et rien que le statut}
\begin{code}
\begin{Verbatim}
$ scripts/check-spec-validated.sh docs/02-specs/SPEC-001.md
KO: docs/02-specs/SPEC-001.md a le statut « Brouillon » (attendu : Validée)

$ scripts/check-spec-validated.sh docs/02-specs/SPEC-001.md
OK: docs/02-specs/SPEC-001.md est Validée
\end{Verbatim}
\end{code}
Le second résultat n'apparaît qu'après qu'un humain a écrit \textbf{Validée} (avec nom et date) dans le tableau d'en-tête.

\begin{retenir}[Avant de proposer « Validée », relisez cette liste]
\begin{itemize}
  \item les sections 1 à 5 sont renseignées ;
  \item aucun marqueur \cmd{[...]} ne subsiste ;
  \item aucune question ouverte bloquante ;
  \item le fichier \cmd{.feature} couvre chaque règle.
\end{itemize}
\end{retenir}

\begin{vousjouez}[À vous de jouer]
Choisissez la première fonctionnalité de votre projet. Écrivez deux règles métier avec leur source. Si l'une n'a pas de source,
écrivez-la en question ouverte et nommez qui peut y répondre. Rédigez trois scénarios (nominal, limite, erreur) en Gherkin français.
\end{vousjouez}

\begin{retenir}[À retenir]
\begin{itemize}
  \item Pas de code sans spec \emph{Validée}, et la validation est humaine.
  \item Règle $\Rightarrow$ source ; règle sans source $\Rightarrow$ question ouverte.
  \item Le script contrôle le statut ; la complétude se relit.
\end{itemize}
\end{retenir}
\source{skills/forge-spec/SKILL.md, templates/docs/02-specs}

\begin{acquis}
  \item Qu'écrit le skill quand une règle métier manque ?
  \item Que contrôle exactement \cmd{check-spec-validated.sh} ?
  \item Combien de scénarios minimum par règle, et de quels types ?
\end{acquis}

\chapter{Décider l'architecture}\label{ch:archi}

\begin{objectifs}[Objectifs --- étape 3, \texttt{forge-architecture}]
\begin{itemize}
  \item reconnaître une décision structurante et la consigner dans un ADR ;
  \item comparer des options \emph{chiffrées} en FCFA, pas des préférences ;
  \item tenir le modèle C4 en code ;
  \item protéger la décision contre la dérive grâce à une fitness function en CI.
\end{itemize}
\end{objectifs}

\begin{situation}[Sur le terrain : \sika{}]
Six mois plus tard, un nouveau développeur demande : \q{Pourquoi n'a-t-on pas fait une application mobile ?}
Si la réponse tient dans \emph{\q{je crois qu'on en avait parlé}}, l'équipe va rediscuter la décision, l'inverser par
inadvertance, ou la refaire pour rien. Un ADR répond en une page, avec les chiffres de l'époque.
\end{situation}

\section{Le concept : décider une fois, par écrit}
\begin{concept}[ADR --- Architecture Decision Record]
Une fiche courte qui consigne \emph{une} décision structurante : le contexte, les options envisagées avec
leur coût, la décision, ses conséquences. Son statut passe de \emph{Proposée} à \emph{Acceptée} sous la
responsabilité d'un humain. Un ADR n'est pas modifié après coup : une décision nouvelle \emph{remplace} l'ancienne
(statut \emph{Remplacée par ADR-XXXX}).
\end{concept}

Sont structurantes : la \textbf{stack technique} (langage et framework du back, du front, du mobile), la base de
données, l'hébergement, toute dépendance majeure, le découpage en modules.

\section{Pas à pas avec \texttt{forge-architecture}}
\begin{enumerate}
  \item Le skill lit les ADR acceptés, le C4 et la context map.
  \item Pour chaque décision, il copie \cmd{adr/0000-template.md} vers \cmd{adr/NNNN-slug.md} et la remplit avec des options chiffrées : coût en FCFA, compétences locales, capacité d'exploitation du client.
  \item Il ne contredit pas l'ADR-0001 (monolithe modulaire) ; en sortir exige un nouvel ADR qui le remplace.
  \item Il tient le C4 en code ; les technologies \cmd{[STACK : à fixer par ADR]} sont remplacées par celles de l'ADR de stack.
  \item Chaque ADR à risque de dérive reçoit une fitness function \cmd{FF-NN}.
\end{enumerate}

\section{L'ADR-0001, fourni : le monolithe modulaire}
Le dossier de conception embarque déjà cette décision, adaptée aux petites équipes et aux clients sans équipe d'exploitation.
Un seul déployable ; \emph{un module par bounded context} ; les modules ne communiquent que par interfaces publiques
ou événements internes ; une base PostgreSQL, un schéma par module.

\begin{figure}[H]
\centering
\resizebox{.95\linewidth}{!}{%
\begin{tikzpicture}[
  box/.style={draw,rounded corners=3pt,align=center,font=\sffamily\small},
  mod/.style={box,fill=accentl!18,draw=accentl,minimum width=2.9cm,minimum height=10mm},
  ext/.style={box,fill=grey!15,draw=grey,minimum width=3cm,minimum height=10mm},
  ar/.style={-{Stealth},thick,grey}]
  \node[box,draw=accent,minimum width=11.4cm,minimum height=4.2cm,fill=accent!4] (api) at (0,0) {};
  \node[font=\sffamily\bfseries\small,text=accent,anchor=north west] at ([xshift=2mm,yshift=-1mm]api.north west) {API --- un seul déployable};
  \node[mod] (c) at (-3.7,0.2) {Cotisation\\{\scriptsize cœur}};
  \node[mod] (n) at (-0.6,0.2) {Notifications};
  \node[mod] (a) at (2.5,0.2) {ACL Paiement};
  \node[font=\sffamily\scriptsize,text=brick] at (-2.15,-1.2) {pas d'import direct entre modules (FF-01)};
  \node[box,fill=codebg,draw=accent!50,minimum width=11cm,minimum height=7mm] (db) at (0,-3.1) {PostgreSQL --- un schéma par module};
  \node[ext] (mm) at (6.8,0.2) {Agrégateur\\Mobile Money};
  \node[ext] (sms) at (1.2,3.1) {Passerelle SMS};
  \node[ext,minimum width=4.2cm] (pwa) at (-3.7,3.1) {PWA (membres, trésorières)};
  \draw[ar] (a)--(mm); \draw[ar] (n.north)--(sms.south west); \draw[ar] (pwa)--(c);
  \draw[ar] (api.south)--(db.north);
\end{tikzpicture}}
\caption{L'ADR-0001 appliqué à \sika{} : un déployable, un module par contexte.}
\end{figure}

\section{L'ADR de stack de \sika{}}
\begin{code}
\begin{Verbatim}
# ADR-0002 — Python (FastAPI) + PostgreSQL pour l'API,
#            PWA légère pour le front

Statut : Acceptée · Décideurs : chef de projet, présidente

## Contexte
Téléphones d'entrée de gamme, data mobile coûteuse,
connexions instables. Pas de DSI ; exploitation par Simetriik.

## Options envisagées (chiffres fictifs)
A. FastAPI + PWA
   + équipe à l'aise, léger   − accès natif limité
   Coût : 25 000 FCFA/mois d'hébergement
B. NestJS + PWA
   + un seul langage          − compétences à acquérir
   Coût : 25 000 FCFA/mois
C. Application native (Flutter)
   + expérience native        − 2 binaires, store, lourd
   Coût : 25 000 FCFA/mois + 150 000 FCFA/an (stores)

## Décision
Option A : une PWA s'installe sans store, charge peu de
données et fonctionne hors ligne pour la consultation.

## Conséquences
- Négatives : pas de notifications push natives avant un
  prochain ADR.
- Fitness function associée : FF-01 (frontières des modules).
\end{Verbatim}
\end{code}

\begin{piege}[Piège : un ADR de préférence]
\q{Je préfère Python} n'est pas une option chiffrée. Un bon ADR répond à la question de la DSI cliente
(ou ici de l'association) : combien cela coûte-t-il, qui peut l'exploiter, et que se passe-t-il si
l'équipe change ?
\end{piege}

\section{Le modèle C4, en code}
\begin{concept}[C4]
Une documentation d'architecture en quatre niveaux de zoom : \emph{contexte}, \emph{conteneurs}, \emph{composants},
\emph{code}. On la décrit en code (Structurizr DSL) plutôt qu'en image : le modèle est versionné, relu en revue et ne
vieillit pas en silence. Les images en sont générées.
\end{concept}
\begin{code}
\begin{Verbatim}
api = container "API (monolithe modulaire)" "Un module par bounded context" "Python / FastAPI" {
    contribution = component "Module Cotisation" "Contexte cœur"
    notification = component "Module Notifications" "Contexte support"
    paymentAcl   = component "ACL Paiement" "Couche anticorruption"
}
paymentAcl -> payment "Initie et vérifie les cotisations" "HTTPS + webhooks"
\end{Verbatim}
\end{code}

\section{Fitness functions : l'architecture sous surveillance}
\begin{concept}[Fitness function]
Un contrôle automatique, exécuté en intégration continue, qui échoue quand une propriété d'architecture est
violée. Chaque ADR à risque de dérive en reçoit une, avec un identifiant \cmd{FF-NN}, un outil et un caractère bloquant ou non.
\end{concept}
\noindent\begin{tabularx}{\linewidth}{llXl}
\toprule
\textbf{Id} & \textbf{Protège} & \textbf{Règle} & \textbf{Bloquant} \\ \midrule
FF-01 & ADR-0001 & aucun import direct entre modules de contextes différents & oui \\
FF-02 & ENF-01 & P95 $<$ 2\,s sur le scénario de référence & non (alerte) \\
FF-03 & ENF-05 & aucune vulnérabilité critique dans les dépendances & oui \\
FF-04 & contrats & l'implémentation respecte \cmd{api/openapi.yaml} & oui \\
FF-05 & glossaire & aucun terme banni dans le code & non \\
\bottomrule
\end{tabularx}
FF-01 et FF-03 se branchent en CI dès le premier lot. Pour \sika{} en Python, FF-01 s'écrit avec \cmd{import-linter} :
\begin{code}
\begin{Verbatim}
[importlinter]
root_package = app

[importlinter:contract:bounded-contexts]
name = Les modules ne s'importent pas entre eux
type = independence
modules =
    app.contribution
    app.notification
    app.payment_acl
\end{Verbatim}
\end{code}
\textbf{Garde-fou :} \cmd{scripts/check-adr-present.sh docs/03-architecture/adr}. Il vérifie qu'au moins un ADR est \emph{Acceptée} ;
il ne lit pas la qualité de l'argumentation.

\begin{vousjouez}[À vous de jouer]
Listez trois décisions de votre projet qu'un nouveau venu pourrait remettre en cause (stack, hébergement, base de données\dots).
Pour la plus structurante, remplissez un ADR avec trois options \emph{chiffrées}. Quelle fitness function protégerait la décision ?
\end{vousjouez}

\begin{retenir}[À retenir]
\begin{itemize}
  \item Une décision structurante = un ADR, avec des options chiffrées ; une nouvelle décision remplace l'ancienne.
  \item Le C4 vit en code, pas en image seule.
  \item Une décision qui peut dériver reçoit une fitness function en CI.
\end{itemize}
\end{retenir}
\source{skills/forge-architecture/SKILL.md, templates/docs/03-architecture}

\begin{acquis}
  \item Citez trois décisions qui exigent un ADR.
  \item Que faut-il pour sortir du monolithe modulaire ?
  \item Quelle fitness function empêche deux contextes de s'importer directement ?
\end{acquis}

\chapter{Le contrat d'API, puis le code}\label{ch:contrat}

\begin{objectifs}[Objectifs --- contrat d'API et implémentation]
\begin{itemize}
  \item écrire une opération OpenAPI rattachée à une spécification validée ;
  \item implémenter en partant d'un test, dans le respect du glossaire et des frontières ;
  \item appliquer les règles FCFA et dates ;
  \item reconnaître les vérifications à passer avant de conclure.
\end{itemize}
\end{objectifs}

\begin{situation}[Sur le terrain : \sika{}]
L'équipe mobile et l'équipe API avancent en parallèle. Sans contrat, la PWA attendra un champ
\cmd{amount} quand l'API enverra \cmd{amountXof} ; l'intégration se fera au dernier jour. Le contrat, écrit
\emph{avant} le code, est la prise de courant que chacun branche de son côté.
\end{situation}

\section{Le contrat précède le code}
\begin{concept}[OpenAPI-first]
Le fichier \cmd{api/openapi.yaml} décrit chaque opération avant son implémentation. Ici, chaque opération \emph{cite}
la spécification dont elle découle : le lien contrat $\leftrightarrow$ spécification est vérifié par un script.
\end{concept}

\subsection*{Les règles du contrat (\texttt{forge-contrat})}
\begin{itemize}
  \item \cmd{operationId} en camelCase anglais, avec les termes du glossaire ;
  \item \cmd{tags} : le bounded context ;
  \item \cmd{summary} citant \cmd{(SPEC-NNN)} --- c'est ce lien que le script lit ;
  \item erreurs en \cmd{application/problem+json} ;
  \item montants : \cmd{integer}, suffixe \cmd{Xof}, \cmd{minimum: 0} ; dates en \cmd{date-time} (UTC) ; identifiants métier avec un \cmd{pattern}.
\end{itemize}
Pour \sika{}, l'opération de la SPEC-001 :
\begin{code}
\begin{Verbatim}
/contributions:
  post:
    operationId: submitContribution
    summary: Cotiser pour le tour en cours (SPEC-001)
    tags: [contribution]
    responses:
      "202": { description: Cotisation initiée }
      "409": { description: Déjà cotisé pour ce tour,
               content: { application/problem+json: {} } }
components:
  schemas:
    Contribution:
      properties:
        reference: { type: string, pattern: "^COT-\d{4}-\d{5}$" }
        amountXof: { type: integer, minimum: 0, example: 5000 }
        initiatedAt: { type: string, format: date-time }
\end{Verbatim}
\end{code}

\subsection*{Le garde-fou \texttt{check-openapi-first.sh}}
Si une opération cite une spec absente ou non validée, le script échoue. Sur \sika{}, \cmd{getPotBalance} cite la SPEC-002 qui n'existe pas encore (sortie réelle) :
\begin{code}
\begin{Verbatim}
$ scripts/check-openapi-first.sh api/openapi.yaml docs/02-specs
KO: getPotBalance → SPEC-002 absente ou non Validée
\end{Verbatim}
\end{code}
Il disparaît lorsque la SPEC-002 est écrite et validée (chapitre~\ref{ch:rapide}).

\begin{concept}[Contrat et test de conformité]
Le contrat sert aussi de base au test FF-04 (Schemathesis) : l'implémentation est comparée automatiquement au contrat.
\skill{forge-contrat} rappelle de le brancher en CI.
\end{concept}

\section{Implémenter avec \texttt{forge-implemente}}
\subsection*{Avant d'écrire une ligne (bloquant)}
\begin{enumerate}[start=0]
  \item \cmd{scripts/check-branch.sh} : une branche \cmd{feat/SPEC-NNN-slug} (jamais \cmd{main}).
  \item \cmd{scripts/check-spec-validated.sh} : la spec est \emph{Validée}, sinon on s'arrête.
  \item Lire le glossaire, les ADR \emph{Acceptée} et la context map.
  \item Le contrat couvre la spec ; sinon on passe par \skill{forge-contrat} d'abord.
  \item Annoncer un plan court et attendre l'accord.
\end{enumerate}

\subsection*{Les règles d'implémentation}
\begin{itemize}
  \item \textbf{Un scénario Gherkin = au moins un test}, écrit \emph{avant} le code.
  \item Code en anglais, commentaires en français, noms du glossaire.
  \item \textbf{FCFA en entiers} (\cmd{amountXof}), jamais de \cmd{float}.
  \item \textbf{Dates} en UTC au stockage, \cmd{Africa/Porto-Novo} à l'affichage.
  \item Aucun import direct entre contextes (FF-01) ; aucune dépendance majeure sans ADR.
  \item Règle absente de la spec : poser la question, ne pas deviner.
\end{itemize}

Voici le test du scénario \q{Montant différent de celui du groupe} (R1 de la SPEC-001), écrit avant la fonction :
\begin{code}
\begin{Verbatim}
def test_contribution_refused_when_amount_differs_from_group():
    """R1 (SPEC-001) : le montant doit être celui fixé pour le groupe."""
    group = SavingsGroup(contribution_amount_xof=5000)
    with pytest.raises(WrongAmount):
        group.accept_contribution(member, amount_xof=3000)
\end{Verbatim}
\end{code}

\begin{piege}[Piège : adapter le test au code]
Un test rouge n'est pas un obstacle à contourner : c'est le scénario qui vous parle. Ne jamais affaiblir un test pour le
faire passer ; si le test est faux, discutez-en d'abord, spec en main.
\end{piege}

\subsection*{Avant de conclure}
Lancer le linter et les tests du projet, puis le garde-fou de vocabulaire :
\cmd{scripts/check-glossary-terms.sh docs/01-domaine/glossaire.md src/}. Un commit porte un seul sujet, un message en anglais au format
Conventional Commits, et cite la spec :
\begin{code}
\begin{Verbatim}
feat(contribution): refuse a contribution whose amount differs from the group

SPEC-001 R1
\end{Verbatim}
\end{code}

\begin{vousjouez}[À vous de jouer]
Pour la spec de votre projet : écrivez l'opération OpenAPI (\cmd{operationId}, \cmd{summary} avec \cmd{(SPEC-NNN)}, une erreur en
\cmd{problem+json}). Lancez \cmd{check-openapi-first.sh}. Écrivez ensuite le test du cas d'erreur \emph{avant} la fonction : à quoi ressemble-t-il ?
\end{vousjouez}

\begin{retenir}[À retenir]
\begin{itemize}
  \item Le contrat se modifie avant le code, et chaque opération cite sa spec.
  \item Un scénario = un test, écrit d'abord.
  \item FCFA entiers, dates UTC, pas d'import entre contextes.
\end{itemize}
\end{retenir}
\source{skills/forge-contrat/SKILL.md, skills/forge-implemente/SKILL.md}

\begin{acquis}
  \item Que doit contenir le \cmd{summary} d'une opération ?
  \item Pourquoi un montant en FCFA n'est-il jamais un \cmd{float} ?
  \item Que faire quand la spec ne dit rien sur un cas rencontré en codant ?
\end{acquis}

\chapter{La recette}\label{ch:recette}

\begin{objectifs}[Objectifs --- étape 5, \texttt{forge-recette}]
\begin{itemize}
  \item dériver le plan de recette des scénarios Gherkin ;
  \item vérifier les critères d'entrée avant de commencer ;
  \item classer les anomalies et gérer leurs délais ;
  \item préparer un procès-verbal que le client, et lui seul, signe.
\end{itemize}
\end{objectifs}

\begin{situation}[Sur le terrain : \sika{}]
Trois trésorières de groupe testent \sika{} sur leur propre téléphone, avec des données anonymisées. L'une d'elles
constate qu'une cotisation confirmée par l'opérateur n'apparaît pas dans le solde du pot. Est-ce grave ?
Réponse du tableau de classification : oui, \emph{bloquante}, car elle empêche le processus métier sans contournement.
\end{situation}

\section{Le concept : la recette est celle du client}
\begin{concept}[Recette]
La phase où le client vérifie, scénario par scénario, que le produit fait ce qui a été validé. Elle ne se confond pas avec
les tests de l'équipe : elle se tient dans un environnement de préproduction, avec des testeurs métier, et se clôt par un
\textbf{procès-verbal (PV)} que le client signe.
\end{concept}

\section{Pas à pas avec \texttt{forge-recette}}
\begin{enumerate}
  \item Le skill lit le plan de recette et chaque \cmd{.feature} des specs du lot.
  \item \textbf{Un scénario Gherkin = une ligne} de la table des scénarios (identifiant, spec, scénario, résultat attendu). Les colonnes \emph{Obtenu}, \emph{Statut}, \emph{Testeur} et \emph{Date} restent vides pour les testeurs.
  \item Il vérifie les \emph{critères d'entrée}.
  \item Il suit les anomalies ; il ne reclasse jamais une anomalie sans l'avis du responsable recette client.
  \item Il prépare le PV ; il ne déclare jamais la recette réussie.
\end{enumerate}

\begin{concept}[Critères d'entrée]
Trois conditions avant de commencer : les specs du lot sont \emph{Validée} ; les scénarios automatisés sont verts en CI ;
un jeu de données anonymisé est chargé. Sinon, la recette ne démarre pas.
\end{concept}

\section{Anomalies et critères de sortie}
\noindent\begin{tabularx}{\linewidth}{llX}
\toprule
\textbf{Niveau} & \textbf{Délai} & \textbf{Définition} \\ \midrule
Bloquante & 48\,h ouvrées & empêche un processus métier sans contournement \\
Majeure & 5 jours ouvrés & processus dégradé, contournement possible \\
Mineure & lot suivant & gêne ergonomique ou cosmétique \\
\bottomrule
\end{tabularx}
\textbf{Critères de sortie :} aucune anomalie bloquante, un nombre borné de majeures avec plan de correction, exigences non fonctionnelles critiques vérifiées.

\subsection*{Un extrait du plan de \sika{}}
\noindent\begin{tabularx}{\linewidth}{lllXl}
\toprule
\textbf{Id} & \textbf{Spec} & \textbf{Scénario} & \textbf{Résultat attendu} & \textbf{Statut} \\ \midrule
SC-01 & SPEC-001 & Cotisation au montant du groupe & cotisation initiée auprès de l'opérateur & OK \\
SC-02 & SPEC-001 & Confirmation par l'opérateur & solde du pot $+$\FCFA{5\,000}, SMS reçu & \textcolor{brick}{KO} \\
SC-03 & SPEC-001 & Seconde cotisation dans le tour & refusée : \q{Déjà cotisé} & OK \\
\bottomrule
\end{tabularx}

\begin{piege}[Piège : corriger à la hâte]
SC-02 échoue. Remplacer l'affichage du solde ferait disparaître le symptôme sans traiter la cause : ici, un webhook de l'opérateur
perdu. Le skill impose la même discipline qu'ailleurs : \emph{reproduire}, poser des hypothèses sur la cause racine, corriger
le minimum via \skill{forge-implemente}, rejouer le scénario. La cause racine retenue : aucune réconciliation périodique des
statuts auprès de l'opérateur. La correction : un rattrapage qui interroge l'opérateur sur les cotisations \q{initiées}.
\end{piege}

\section{Le procès-verbal}
\begin{code}
\begin{Verbatim}
Lot : 1 · Date : [JJ/MM/AAAA] · Décision : [ ] Acceptée [ ] Acceptée avec réserves [ ] Refusée
Réserves :
Signatures : Client ____________ Simetriik Solutions ____________
\end{Verbatim}
\end{code}
Le PV est la porte d'entrée de la production (chapitre~\ref{ch:prod}). Aucune tâche, aucun assistant ne signe à la place du client.

\begin{vousjouez}[À vous de jouer]
Prenez les scénarios de la spec de votre projet et transformez-les en lignes de plan de recette. Qui, chez votre client, jouera le
rôle de responsable recette ? Et quelle anomalie classeriez-vous \emph{bloquante} dans votre contexte ? Écrivez-la.
\end{vousjouez}

\begin{retenir}[À retenir]
\begin{itemize}
  \item Un scénario Gherkin = une ligne du plan de recette.
  \item Critères d'entrée avant de démarrer ; critères de sortie avant de conclure.
  \item Le client classe les anomalies ; le client signe le PV.
\end{itemize}
\end{retenir}
\source{skills/forge-recette/SKILL.md, templates/docs/05-recette}

\begin{acquis}
  \item Citez les trois critères d'entrée d'une recette.
  \item Qui peut reclasser une anomalie ?
  \item Que fait-on d'une anomalie bloquante découverte en recette ?
\end{acquis}

\chapter{Mettre en production}\label{ch:prod}

\begin{objectifs}[Objectifs --- étape 6, \texttt{forge-prod}]
\begin{itemize}
  \item vérifier la porte d'entrée : PV signé, aucune anomalie bloquante ;
  \item rédiger plan de déploiement, runbook, retour arrière et monitoring ;
  \item relier chaque objectif de l'impact map à un indicateur suivi après la mise en service ;
  \item savoir quelles actions exigent votre confirmation explicite.
\end{itemize}
\end{objectifs}

\begin{situation}[Sur le terrain : \sika{}]
Vendredi, 17 h. Le PV du lot 1 est signé. Un coéquipier propose : \q{On met en prod ce soir, ça passera.}
Or le marché ferme à 18 h, les cotisations de la semaine tombent entre 16 h et 20 h, et personne n'a dit
comment revenir en arrière si la confirmation des opérateurs échoue. Cette étape existe pour que la
mise en service soit un non-événement.
\end{situation}

\section{La porte d'entrée}
Le skill commence par une vérification qui peut tout arrêter : le \textbf{PV de recette est-il signé par le client, sans
anomalie bloquante ?} Sinon, il s'arrête et le dit.

\section{Pas à pas avec \texttt{forge-prod}}
Il produit quatre documents dans \cmd{docs/06-production/} :

\noindent\begin{tabularx}{\linewidth}{lX}
\toprule
\textbf{Livrable} & \textbf{Contenu pour \sika{}} \\ \midrule
\cmd{plan-deploiement.md} & hébergement (une VM, Docker Compose, d'après l'ADR-0001), pipeline CI/CD, fenêtre de déploiement \\
\cmd{runbook.md} & que faire si les confirmations de l'opérateur n'arrivent plus ? comment relancer la réconciliation ? \\
\cmd{rollback.md} & retour arrière avec déclencheurs \emph{chiffrés} \\
\cmd{plan-monitoring.md} & indicateurs techniques et indicateurs métier liés aux objectifs \\
\bottomrule
\end{tabularx}

Le skill interroge d'abord sur l'hébergement réel (VM du client, cloud, on-premise), la capacité d'exploitation et le budget
en FCFA. Une décision d'hébergement ou d'orchestration est un ADR, via \skill{forge-architecture}.

\subsection*{Les contraintes par défaut}
\begin{itemize}
  \item connexions lentes ou instables : le service doit tolérer coupures et reprises ;
  \item paiements Mobile Money : une file de reprise, des statuts \emph{vérifiables auprès de l'opérateur} ;
  \item secrets en variables d'environnement, jamais dans le code.
\end{itemize}

\subsection*{Retour arrière : des déclencheurs, pas des intentions}
\begin{code}
\begin{Verbatim}
Déclencheurs de retour arrière (Sika, lot 1)
- plus de 5 % de cotisations non confirmées au bout de 15 min, pendant 15 min
- aucune confirmation reçue d'un opérateur pendant 10 min
- erreur 5xx sur plus de 2 % des requêtes pendant 10 min
\end{Verbatim}
\end{code}

\subsection*{Monitoring métier : relier l'impact map à la production}
Chaque objectif \texttt{O\#} reçoit un indicateur, une source de mesure et deux revues.
\noindent\begin{tabularx}{\linewidth}{llXll}
\toprule
\textbf{Objectif} & \textbf{Indicateur} & \textbf{Source} & \textbf{J+7} & \textbf{J+30} \\ \midrule
O1 & \% de cotisations contestées & journal des contestations, relevé de la trésorière & revue & revue \\
O2 & délai de remise du pot & horodatages \emph{Tour clôturé} $\to$ \emph{Pot remis} & première mesure & revue \\
\bottomrule
\end{tabularx}
La valeur \emph{À mesurer} de O2 trouve enfin sa réponse : le journal d'événements de \sika{} la produira.

\begin{piege}[Piège : « le skill l'a déployé »]
Ce skill \emph{rédige et prépare} ; il n'exécute \textbf{jamais} un déploiement de production. Toute modification de CI/CD, de Docker, de
fichiers \cmd{.env}, toute migration, toute donnée de production, tout \cmd{git push} et tout appel à un service payant exigent votre
confirmation explicite.
\end{piege}

\section{La clôture}
Après la mise en service : revue post-production à J+7 et J+30, puis retour à \skill{forge-cadrage} pour le lot suivant. La boucle est
fermée : les mesures de production nourrissent le prochain cadrage.

\begin{vousjouez}[À vous de jouer]
Pour votre projet, écrivez trois déclencheurs chiffrés de retour arrière. Puis choisissez un objectif de votre impact map : où
sera-t-il mesuré, par qui, à J+7 et J+30 ?
\end{vousjouez}

\begin{retenir}[À retenir]
\begin{itemize}
  \item Porte d'entrée : PV signé, zéro bloquante.
  \item Quatre livrables : déploiement, runbook, retour arrière, monitoring.
  \item Chaque objectif $\to$ un indicateur $\to$ une revue à J+7 et J+30.
  \item Le skill prépare ; vous confirmez toute action sensible.
\end{itemize}
\end{retenir}
\source{skills/forge-prod/SKILL.md, templates/docs/06-production}

\begin{acquis}
  \item Quelle est la porte d'entrée de la production ?
  \item Pourquoi un déclencheur de retour arrière doit-il être chiffré ?
  \item Quelles actions exigent toujours votre confirmation ?
\end{acquis}

\part{Au quotidien}
\chapter{Le cycle court}\label{ch:rapide}

\begin{objectifs}[Objectifs --- \texttt{forge-rapide}]
\begin{itemize}
  \item décider si une demande relève du cycle court ou du cycle complet ;
  \item rédiger une mini-spec d'une seule règle ;
  \item comprendre le garde-fou qui empêche le cycle court de devenir une porte dérobée.
\end{itemize}
\end{objectifs}

\begin{situation}[Sur le terrain : \sika{}]
Une trésorière demande : \q{Je voudrais voir le montant total du pot sur l'écran du tour.} Parcourir sept étapes pour
afficher une somme serait absurde. Mais afficher une somme n'est pas \q{juste} un écran : quelle somme ? les
cotisations initiées comptent-elles ? Le cycle court garde la rigueur et allège la cérémonie.
\end{situation}

\section{Quand l'utiliser : six questions}
\begin{figure}[H]
\centering
\resizebox{\linewidth}{!}{%
\begin{tikzpicture}[
  q/.style={draw=accent,fill=accent!6,rounded corners=3pt,align=center,font=\sffamily\small,text width=11cm,minimum height=8mm},
  no/.style={draw=leaf,fill=leaf!12,rounded corners=3pt,align=center,font=\sffamily\small\bfseries,minimum height=9mm,text width=11cm},
  yes/.style={draw=brick,fill=brick!10,rounded corners=3pt,align=center,font=\sffamily\small\bfseries,text width=4.2cm,minimum height=26mm},
  ar/.style={-{Stealth},thick,grey}]
  \node[q] (q) at (0,0) {Nouveau bounded context ? Nouvelle dépendance ? Migration de données ?\\Rupture du contrat OpenAPI ? Plus d'une règle métier ? Décision structurante (ADR) ?};
  \node[no] (r) at (0,-2.2) {Six réponses « non » $\Rightarrow$ cycle court (\skill{forge-rapide})};
  \node[yes] (y) at (8.3,-1) {Un seul « oui » $\Rightarrow$ cycle complet (\skill{forge-spec})};
  \draw[ar] (q)--(r) node[midway,right,font=\sffamily\scriptsize]{tous non};
  \draw[ar] (q.east)--++(.6,0)|-(y.west) node[pos=.2,above,font=\sffamily\scriptsize]{au moins un oui};
\end{tikzpicture}}
\caption{L'éligibilité au cycle court.}
\end{figure}

\section{Déroulement}
\begin{enumerate}
  \item \textbf{Mini-spec.} Copiez \cmd{_TEMPLATE-spec-rapide.md} en \cmd{SPEC-NNN.md} : une règle (avec sa source), un scénario, l'\cmd{operationId} si l'API change, taille \texttt{S}. Seul un humain la passe à \emph{Validée}.
  \item \textbf{Réalisation.} Branche \cmd{feat/SPEC-NNN-slug}, contrat si l'API change, test d'abord, puis code.
  \item \textbf{Vérifications.} Linter, tests, \cmd{check-glossary-terms.sh}, et \cmd{check-rapide-eligible.sh}.
\end{enumerate}

\subsection*{La mini-spec de \sika{}}
\begin{code}
\begin{Verbatim}
# SPEC-002 — Afficher le solde du pot

Statut : Validée · Taille : S · Contexte : Cotisation
Objectif lié : O1

## Règle
Le solde du pot est la somme des cotisations confirmées du tour.
Source : atelier Event Storming, politique « dès qu'une
cotisation est confirmée ».

## Scénario
Scénario: Solde du pot
  Étant donné deux cotisations confirmées de 5000 FCFA
  Et une cotisation initiée
  Quand la trésorière consulte le pot
  Alors le solde affiché est de 10000 FCFA

## Interface (si l'API change)
operationId : getPotBalance
\end{Verbatim}
\end{code}
La règle est la même que R2 de la SPEC-001, vue autrement : le solde ne compte que les cotisations \emph{confirmées}. Le cas du
doute (initiée, non confirmée) est dans le scénario lui-même.

\section{Le garde-fou d'éligibilité}
Pour que le cycle court ne devienne pas une porte dérobée, \cmd{scripts/check-rapide-eligible.sh docs/02-specs/SPEC-NNN.md} échoue si :
\begin{itemize}
  \item la spec n'est pas de taille \texttt{S} ;
  \item plus de 8 fichiers sont modifiés par rapport à \cmd{main} ;
  \item un manifeste de dépendances change (\cmd{package.json}, \cmd{pyproject.toml}, \cmd{go.mod}\dots) ;
  \item une migration ou un fichier \cmd{.sql} apparaît.
\end{itemize}
En cas d'échec, le travail a dépassé le cadre : basculez vers le cycle complet.

\begin{piege}[Piège : le faux refus vaut mieux que l'oubli]
Une modification de \cmd{pyproject.toml} est refusée même sans nouvelle dépendance (une simple version de l'outil, par
exemple). Le script préfère un faux refus à une dépendance oubliée. Le seuil de 8 fichiers est arbitraire : réglez-le à l'usage.
\end{piege}

\begin{vousjouez}[À vous de jouer]
Prenez trois demandes récentes de votre projet. Pour chacune, répondez aux six questions : laquelle aurait pu être traitée en
cycle court ? Rédigez la mini-spec de celle-là.
\end{vousjouez}

\begin{retenir}[À retenir]
\begin{itemize}
  \item Cycle court = une règle, taille \texttt{S}, ni dépendance ni migration ni rupture de contrat.
  \item La mini-spec reste obligatoire et validée par un humain.
  \item Le script d'éligibilité vous ramène au cycle complet quand il le faut.
\end{itemize}
\end{retenir}
\source{skills/forge-rapide/SKILL.md, scripts/check-rapide-eligible.sh}

\begin{acquis}
  \item Citez trois raisons de refuser le cycle court.
  \item La mini-spec peut-elle être validée par l'assistant ?
  \item Que signifie un échec de \cmd{check-rapide-eligible.sh} ?
\end{acquis}

\chapter{Rigueur Git}\label{ch:git}

\begin{objectifs}[Objectifs]
\begin{itemize}
  \item nommer vos branches et vos commits selon les conventions du projet ;
  \item installer les hooks Git et comprendre ce qu'ils bloquent ;
  \item connaître les limites des hooks locaux.
\end{itemize}
\end{objectifs}

\begin{situation}[Sur le terrain : \sika{}]
Un vendredi soir, un correctif \q{rapide} est poussé directement sur \cmd{main}, sans spec. Lundi, la trésorière générale
signale un comportement que personne ne sait expliquer. La rigueur Git n'est pas de la bureaucratie : c'est la
traçabilité de qui a changé quoi, pourquoi, et avec quelle autorisation.
\end{situation}

\section{Les règles}
\begin{itemize}
  \item \textbf{Une branche par fonctionnalité ou correction} : \cmd{feat/SPEC-NNN-slug} (par exemple \cmd{feat/SPEC-001-cotiser}) ou \cmd{fix/slug}. Jamais de code sur \cmd{main}.
  \item \textbf{Commits petits et atomiques}, un sujet par commit, messages en anglais au format Conventional Commits : \cmd{feat}, \cmd{fix}, \cmd{refactor}, \cmd{docs}, \cmd{test}, \cmd{chore}.
  \item Le corps du commit cite \cmd{SPEC-NNN}.
  \item Aucun secret, aucun fichier généré lourd, aucun \cmd{node_modules}.
\end{itemize}

\section{Les scripts, en action}
\noindent\begin{tabularx}{\linewidth}{>{\raggedright\arraybackslash}p{5.2cm}X}
\toprule
\textbf{Script} & \textbf{Rôle} \\ \midrule
\cmd{check-branch.sh} & refuse \cmd{main}, exige \cmd{feat/...} ou \cmd{fix/...} \\
\cmd{check-commit-msg.sh} & Conventional Commits, 72 caractères maximum sur la première ligne \\
\cmd{check-commit-allowed.sh} & sur \cmd{feat/SPEC-NNN-slug}, refuse le code tant que la spec n'est pas \emph{Validée} \\
\cmd{install-hooks.sh} & installe les hooks \cmd{pre-commit} et \cmd{commit-msg} \\
\bottomrule
\end{tabularx}
\begin{code}
\begin{Verbatim}
$ scripts/check-branch.sh feat/SPEC-001-cotiser
OK: branche feat/SPEC-001-cotiser

$ scripts/check-commit-msg.sh msg-bad.txt      # « added stuff »
KO: 'added stuff' — attendu « type(portée): sujet » (feat|fix|refactor|docs|test|chore), 72 car. max

$ scripts/check-commit-msg.sh msg-ok.txt       # « feat(api): add login »
OK: message conforme
\end{Verbatim}
\end{code}

\section{Le hook \emph{pre-commit}}
\begin{figure}[H]
\centering
\resizebox{\linewidth}{!}{%
\begin{tikzpicture}[
  s/.style={draw=accent,fill=accent!6,rounded corners=3pt,align=center,font=\sffamily\small,minimum width=3.3cm,minimum height=10mm},
  ok/.style={s,draw=leaf,fill=leaf!12},
  ko/.style={s,draw=brick,fill=brick!10},
  ar/.style={-{Stealth},thick,grey},
  l/.style={font=\sffamily\scriptsize,text=grey}]
  \node[s] (a) at (0,0) {git commit};
  \node[s] (b) at (4.4,0) {Seulement\\docs, .md, tests ?};
  \node[ok] (o1) at (4.4,-2.2) {commit accepté};
  \node[s] (c) at (8.8,0) {Branche\\fix/... ?};
  \node[s] (d) at (13.2,0) {Spec nommée\\Validée ?};
  \node[ko] (k) at (13.2,-2.2) {commit refusé};
  \node[ok] (o2) at (8.8,-2.2) {commit accepté};
  \draw[ar] (a)--(b); \draw[ar] (b)--(c) node[midway,above,l]{non}; \draw[ar] (b)--(o1) node[midway,right,l]{oui};
  \draw[ar] (c)--(d) node[midway,above,l]{non}; \draw[ar] (c)--(o2) node[midway,right,l]{oui};
  \draw[ar] (d)--(k) node[midway,right,l]{non}; \draw[ar] (d.south west)--(o2.north east) node[midway,above,l,sloped]{oui};
\end{tikzpicture}}
\caption{Les décisions du hook \emph{pre-commit}.}
\end{figure}
Installation, sur confirmation (le script modifie le dépôt) : \cmd{scripts/install-hooks.sh}. Il n'écrase jamais un hook existant.
Une branche \cmd{feat/} sans \cmd{SPEC-NNN} est refusée avec le nommage attendu.

\begin{piege}[Piège : contourner et oublier]
Avec les hooks installés, tout commit sur \cmd{main} est refusé, y compris de documentation. Et \cmd{git commit --no-verify}
contourne le hook : seule une CI côté serveur rendrait la règle absolue. Si vous contournez, c'est un choix assumé et dit.
\end{piege}

\section{Les pull requests}
Le modèle \cmd{.github/pull_request_template.md} demande la ou les specs concernées, l'ADR concerné, le lot ou la version, et la liste des
vérifications : spec validée, glossaire respecté, contrat à jour, un test par critère d'acceptation, fitness functions, branche et messages conformes.

\begin{vousjouez}[À vous de jouer]
Installez les hooks sur un dépôt d'essai. Essayez de committer un fichier de code sur \cmd{main}, puis sur
\cmd{feat/SPEC-001-essai} avec une spec en \emph{Brouillon}. Notez chaque message de refus : à quelle règle correspond-il ?
\end{vousjouez}

\begin{retenir}[À retenir]
\begin{itemize}
  \item Une branche par sujet ; jamais de code sur \cmd{main}.
  \item Conventional Commits en anglais, citant \cmd{SPEC-NNN}.
  \item Les hooks sont locaux et contournables : la CI serveur est le filet final.
\end{itemize}
\end{retenir}
\source{scripts/check-branch.sh, check-commit-msg.sh, check-commit-allowed.sh, install-hooks.sh}

\begin{acquis}
  \item Quels commits le hook \emph{pre-commit} laisse-t-il toujours passer ?
  \item Pourquoi \cmd{--no-verify} limite-t-il la portée des hooks ?
  \item Dans quelle langue écrit-on un message de commit ?
\end{acquis}

\chapter{S'entraîner avec \texttt{forge-apprendre}}\label{ch:apprendre}

\begin{objectifs}[Objectifs]
\begin{itemize}
  \item suivre quatre leçons pratiques sur un cas fictif, avec un coach ;
  \item mesurer votre progression sans fichier de suivi ;
  \item transférer chaque leçon sur votre projet réel.
\end{itemize}
\end{objectifs}

\begin{situation}[Sur le terrain : une nouvelle recrue]
Une ingénieure rejoint l'équipe de \sika{}. Elle ne veut pas lire un manuel de bout en bout : elle veut essayer. Le skill
\skill{forge-apprendre} lui donne un coach et un terrain d'essai.
\end{situation}

\section{Un coach, pas un exécutant}
Le skill \emph{explique}, \emph{relit} et donne des indices à la demande ; c'est \textbf{vous} qui écrivez chaque
livrable. Les exercices portent sur un autre cas fictif que \sika{}, \textbf{PayBoutique} (une boutique en ligne béninoise
qui encaisse des paiements MTN MoMo et Moov Money), et vivent dans \cmd{apprentissage/}, jamais dans \cmd{docs/} : le vrai dossier
de conception reste intact.

\begin{concept}[Apprentissage par la pratique]
Chaque leçon suit le même schéma : le \emph{concept} en quelques lignes, le \emph{pourquoi}, puis un \emph{exercice} vérifié par un garde-fou. À la fin,
vous refaites l'exercice sur votre projet avec le skill de l'étape.
\end{concept}

\section{Les quatre leçons}
\noindent\begin{tabularx}{\linewidth}{lXX}
\toprule
\textbf{Leçon} & \textbf{Exercice} & \textbf{Vérification} \\ \midrule
1 Impact map & objectif mesurable, acteur, tri de trois idées & mots-clés \emph{indicateur}, \emph{cible}, \emph{échéance} \\
2 Langage & glossaire et code contenant un terme banni & \cmd{check-glossary-terms.sh} : KO puis OK \\
3 Spec & règle de paiement avec trois cas & \cmd{check-spec-validated.sh} \\
4 ADR & choix d'hébergement chiffré en FCFA & \cmd{check-adr-present.sh} \\
\bottomrule
\end{tabularx}

\section{Suivre sa progression}
\cmd{scripts/forge-lecons.sh} affiche par exemple \cmd{Leçons [██░░] 2/4 · prochaine : 03-spec}, calculé à partir de vos fichiers d'exercice : vous
reprenez là où vous vous étiez arrêté, sans fichier de suivi.

\begin{piege}[Piège : croire que le vert prouve la qualité]
Les scripts ne vérifient que la présence et le statut, pas la qualité. La qualité se discute en relecture avec le coach. Les concepts
sans leçon (bounded context, OpenAPI d'abord, fitness functions, recette) s'étudient à partir du skill correspondant.
\end{piege}

\begin{vousjouez}[À vous de jouer]
Lancez \skill{forge-apprendre}, faites la leçon~1, puis refaites immédiatement l'exercice sur votre projet avec \skill{forge-cadrage}.
Qu'est-ce qui était différent quand les chiffres étaient les vôtres ?
\end{vousjouez}

\begin{retenir}[À retenir]
\begin{itemize}
  \item Le coach explique ; vous écrivez.
  \item Quatre leçons : impact map, langage, spec, ADR.
  \item Après chaque leçon, transférez sur votre projet avec le skill de l'étape.
\end{itemize}
\end{retenir}
\source{skills/forge-apprendre/SKILL.md et lecons/}

\begin{acquis}
  \item Où vivent les exercices, et pourquoi pas dans \cmd{docs/} ?
  \item Que ne vérifie \emph{pas} \cmd{forge-lecons.sh} ?
  \item Quelle leçon se rattache à quel skill d'étape ?
\end{acquis}

\chapter{La chaîne de traçabilité}\label{ch:trace}

\begin{objectifs}[Objectifs]
\begin{itemize}
  \item suivre un objectif jusqu'à son indicateur de production, et inversement ;
  \item repérer à quelle étape une rupture de la chaîne se détecte ;
  \item répondre à la question : « pourquoi ce code existe-t-il ? »
\end{itemize}
\end{objectifs}

Tout ce manuel converge vers une idée : \emph{chaque ligne de code doit pouvoir justifier son existence}. Voici, pour \sika{}, le chemin complet
d'un objectif, de la première réunion à la mesure en production.

\begin{figure}[H]
\centering
\resizebox{\linewidth}{!}{%
\begin{tikzpicture}[
  n/.style={draw,rounded corners=3pt,align=center,font=\sffamily\small,minimum width=3.3cm,minimum height=13mm},
  c0/.style={n,fill=accent!8,draw=accent},
  c1/.style={n,fill=ochre!12,draw=ochre},
  c2/.style={n,fill=leaf!12,draw=leaf},
  ar/.style={-{Stealth},thick,grey},
  l/.style={font=\sffamily\scriptsize,text=grey,align=center}]
  \node[c0] (o) at (0,0) {\textbf{O1}\\cotisations contestées};
  \node[c0] (f) at (4.4,0) {\textbf{Fonctionnalité}\\Cotiser par MoMo};
  \node[c1] (s) at (8.8,0) {\textbf{SPEC-001}\\règle R2};
  \node[c1] (g) at (13.2,0) {\textbf{.feature}\\scénario \q{Confirmation}};
  \node[c2] (a) at (13.2,-2.6) {\textbf{openapi.yaml}\\submitContribution};
  \node[c2] (t) at (8.8,-2.6) {\textbf{Test}\\\texttt{test\_confirm...}};
  \node[c2] (r) at (4.4,-2.6) {\textbf{Recette}\\SC-02};
  \node[c0] (m) at (0,-2.6) {\textbf{Production}\\\% contestées à J+30};
  \draw[ar] (o)--(f); \draw[ar] (f)--(s); \draw[ar] (s)--(g); \draw[ar] (g)--(a); \draw[ar] (a)--(t); \draw[ar] (t)--(r); \draw[ar] (r)--(m);
  \draw[ar,dashed,accent] (m.north)--(o.south) node[midway,right,l]{boucle de retour};
\end{tikzpicture}}
\caption{La chaîne de traçabilité d'un objectif de \sika{} : chaque maillon est contrôlé par un artefact ou un garde-fou.}
\end{figure}

\section{Où une rupture se détecte}
\noindent\begin{tabularx}{\linewidth}{>{\raggedright\arraybackslash}p{4.6cm}X}
\toprule
\textbf{Rupture} & \textbf{Détectée par} \\ \midrule
Fonctionnalité sans objectif & colonne « Objectif » de l'impact map, \skill{forge-cadrage} \\
Terme utilisé sous deux noms & \cmd{check-glossary-terms.sh} (FF-05) \\
Code sans spec validée & hook \emph{pre-commit}, \cmd{check-spec-validated.sh} \\
Opération d'API sans spec & \cmd{check-openapi-first.sh} \\
Décision sans trace & \cmd{check-adr-present.sh}, ADR \\
Contexte qui en importe un autre & FF-01 en CI \\
Travail qui déborde du cycle court & \cmd{check-rapide-eligible.sh} \\
Production sans mesure & plan de monitoring, revues J+7 et J+30 \\
\bottomrule
\end{tabularx}

\section{Votre carnet de bord, version finale}
\begin{vousjouez}[À vous de jouer : votre propre chaîne]
Choisissez \emph{une} fonctionnalité de votre projet. Écrivez, à la main, les huit maillons de la figure ci-dessus avec les identifiants réels : objectif,
spec, règle, scénario, opération, test, scénario de recette, indicateur. S'il vous manque un maillon, vous venez de trouver la prochaine tâche du projet.
\end{vousjouez}

\begin{retenir}[À retenir]
\begin{itemize}
  \item Chaque ligne de code remonte à une règle, une spec, une fonctionnalité, un objectif.
  \item Les garde-fous contrôlent les maillons mécaniques ; la relecture humaine contrôle le sens.
  \item La boucle se ferme : la mesure en production nourrit le cadrage du lot suivant.
\end{itemize}
\end{retenir}


\appendix
\addpart{Annexes}
\chapter{Référence des commandes et scripts}
\section{Skills}
\noindent\begin{tabularx}{\linewidth}{>{\raggedright\arraybackslash}p{5.6cm}X}
\toprule \textbf{Skill} & \textbf{Rôle} \\ \midrule
\cmd{forge-init} & instancie le dossier de conception \\
\cmd{forge-cadrage} & étape 0 : vision et impact map \\
\cmd{forge-domaine} & étape 1 : event storming, glossaire, context map \\
\cmd{forge-spec} & étape 2 : spécification et Gherkin \\
\cmd{forge-architecture} & étape 3 : C4, ADR (dont la stack), fitness functions \\
\cmd{forge-contrat} & contrat OpenAPI rattaché aux specs \\
\cmd{forge-implemente} & implémentation d'une spec validée \\
\cmd{forge-rapide} & cycle court pour petites fonctionnalités \\
\cmd{forge-recette} & étape 5 : plan de recette \\
\cmd{forge-prod} & étape 6 : déploiement, runbook, rollback, monitoring \\
\cmd{forge-statut} & tableau de bord et prochaine action \\
\cmd{forge-apprendre} & leçons pratiques \\
\bottomrule
\end{tabularx}

\section{Scripts (\texttt{scripts/})}
\noindent\begin{tabularx}{\linewidth}{>{\raggedright\arraybackslash}p{5.6cm}X}
\toprule \textbf{Script} & \textbf{Contrôle} \\ \midrule
\cmd{check-spec-validated.sh} & une spec est au statut \emph{Validée} (statut seulement) \\
\cmd{check-openapi-first.sh} & chaque \cmd{operationId} cite une spec validée \\
\cmd{check-adr-present.sh} & au moins un ADR \emph{Acceptée} \\
\cmd{check-glossary-terms.sh} & aucun synonyme banni dans le code (FF-05) \\
\cmd{check-branch.sh} & nom de branche conforme, pas \cmd{main} \\
\cmd{check-commit-msg.sh} & message Conventional Commits \\
\cmd{check-commit-allowed.sh} & spec validée avant tout commit de code \\
\cmd{check-rapide-eligible.sh} & le travail reste éligible au cycle court \\
\cmd{install-hooks.sh} & installe les hooks Git \\
\cmd{forge-status.sh} & état de la démarche \\
\cmd{forge-progress.sh} & barre de progression 0 à 6 \\
\cmd{forge-lecons.sh} & progression des leçons \\
\bottomrule
\end{tabularx}
Tous sont du bash pur (avec awk et grep) et réutilisables en CI. \cmd{bash tests/run.sh} exécute leurs tests.

\chapter{Glossaire}\label{app:glossaire}
\begin{description}[style=nextline]
  \item[ADR] \emph{Architecture Decision Record} : fiche qui consigne une décision structurante (contexte, options chiffrées, décision, conséquences).
  \item[AGENTS.md] Fichier de consignes relu par l'assistant de code à chaque session.
  \item[Bounded context] Frontière à l'intérieur de laquelle un modèle et son vocabulaire sont cohérents.
  \item[C4] Documentation d'architecture en quatre niveaux : contexte, conteneurs, composants, code.
  \item[Context map] Carte des bounded contexts et de leurs relations.
  \item[Cycle court] Parcours allégé pour une petite fonctionnalité (une règle, taille S), contrôlé par un garde-fou d'éligibilité.
  \item[Event Storming] Atelier de modélisation du domaine à partir des événements métier.
  \item[Fitness function] Contrôle automatique, exécuté en CI, qui protège une propriété d'architecture (FF-01 à FF-05).
  \item[Garde-fou] Script déterministe qui échoue quand une règle de la démarche est violée.
  \item[Gherkin] Langage d'exemples (Étant donné / Quand / Alors) qui sert de critère d'acceptation et de test.
  \item[Impact map] Chaîne objectif $\rightarrow$ acteur $\rightarrow$ comportement $\rightarrow$ fonctionnalité.
  \item[Langage omniprésent] Vocabulaire unique partagé par le métier et le code (glossaire).
  \item[MoSCoW] Priorisation : Must, Should, Could, Won't.
  \item[Objectif mesurable] Objectif avec indicateur, valeur actuelle, cible et échéance.
  \item[OpenAPI] Description normalisée d'une API ; ici, écrite \emph{avant} le code.
  \item[PV de recette] Procès-verbal signé par le client qui clôt la recette.
  \item[Skill] Instruction de travail que l'assistant de code suit pour une étape.
\end{description}

\chapter{Dépannage}\label{app:depannage}
\begin{description}[style=nextline]
  \item[Les skills n'apparaissent pas] Sous Claude Code, vérifier \cmd{--plugin-dir}. Sous Codex ou OpenCode, relancer \cmd{./install.sh} puis redémarrer l'outil.
  \item[\texttt{forge-init} s'arrête] Un dossier \cmd{docs/} ou un fichier \cmd{AGENTS.md} existe déjà : le skill n'écrase jamais.
  \item[Commit refusé : branche] Créez une branche \cmd{feat/SPEC-NNN-slug} ou \cmd{fix/slug} ; les hooks interdisent \cmd{main}.
  \item[Commit refusé : spec non validée] La spec \cmd{SPEC-NNN} nommée dans la branche n'est pas \emph{Validée}. Faites-la valider par un humain ou, pour un commit de documentation seule, ne touchez que \cmd{docs/} et les fichiers \cmd{.md}.
  \item[Commit refusé : message] Le message doit suivre \cmd{type(portée): sujet}, 72 caractères au plus sur la première ligne.
  \item[Cycle court refusé] \cmd{check-rapide-eligible.sh} a détecté un dépassement (taille, fichiers, dépendances, migration) : passez par \skill{forge-spec}.
  \item[La barre reste à 0/7] Le livrable n'est pas au statut \emph{Validée} (ADR : \emph{Acceptée}) ; vérifier la ligne \cmd{Statut} en tête du fichier.
  \item[Terme banni détecté] Remplacez le synonyme par le nom technique du glossaire ; n'affaiblissez pas le glossaire pour faire passer le contrôle.
  \item[Le contrat est refusé] Une opération cite une spec absente ou non validée : validez la spec, ou corrigez le \cmd{summary}.
\end{description}

\chapter{Corrigés des auto-évaluations}\label{app:corriges}
Comparez vos réponses, puis relisez le chapitre indiqué si un point vous échappe.

\section*{Chapitre~\ref{ch:demarche} --- Pourquoi une démarche}
\begin{enumerate}[label=\sffamily\bfseries\color{leaf}Q\arabic*.,leftmargin=2.2em]
  \item Un skill \emph{guide} l'assistant (questions, gabarits) ; un garde-fou \emph{contrôle} un fait et échoue si la règle est violée.
  \item Un humain, avec son nom et la date. Jamais l'assistant.
  \item Le contrat d'API d'abord ; le code ensuite.
\end{enumerate}
\section*{Chapitre~\ref{ch:prise} --- Prise en main}
\begin{enumerate}[label=\sffamily\bfseries\color{leaf}Q\arabic*.,leftmargin=2.2em]
  \item Le skill s'arrête et le signale : il n'écrase jamais.
  \item Choisir la stack est une décision d'architecture (étape~3, ADR), à prendre après le cadrage et le domaine.
  \item Le livrable doit être \emph{Validée} (ADR : \emph{Acceptée}) ; pour l'architecture, au moins deux ADR acceptés.
\end{enumerate}
\section*{Chapitre~\ref{ch:cadrage} --- Cadrage}
\begin{enumerate}[label=\sffamily\bfseries\color{leaf}Q\arabic*.,leftmargin=2.2em]
  \item Il n'a ni indicateur, ni valeur actuelle, ni cible, ni échéance : on ne saura jamais s'il est atteint.
  \item Il la signale hors périmètre.
  \item Le sponsor client.
\end{enumerate}
\section*{Chapitre~\ref{ch:domaine} --- Domaine}
\begin{enumerate}[label=\sffamily\bfseries\color{leaf}Q\arabic*.,leftmargin=2.2em]
  \item Il décrit un fait déjà survenu (\q{Cotisation confirmée}), ce qui oblige à nommer ce qui s'est réellement passé.
  \item Les synonymes à proscrire ; le script \cmd{check-glossary-terms.sh} les cherche dans le code.
  \item Pour que le cœur ne parle jamais le langage ni les particularités d'un opérateur Mobile Money.
\end{enumerate}
\section*{Chapitre~\ref{ch:spec} --- Spécification}
\begin{enumerate}[label=\sffamily\bfseries\color{leaf}Q\arabic*.,leftmargin=2.2em]
  \item Il l'inscrit en \emph{Questions ouvertes} et vous la demande.
  \item Le statut \emph{Validée} et la présence d'un nom et d'une date sur chaque ligne \cmd{Validé par}. Ni la complétude ni la qualité : elles se relisent à la main.
  \item Au moins trois : un nominal, un limite, un d'erreur.
\end{enumerate}
\section*{Chapitre~\ref{ch:archi} --- Architecture}
\begin{enumerate}[label=\sffamily\bfseries\color{leaf}Q\arabic*.,leftmargin=2.2em]
  \item Par exemple : la stack, l'hébergement, la base de données, une dépendance majeure, le découpage en modules.
  \item Un nouvel ADR qui remplace l'ADR-0001.
  \item FF-01.
\end{enumerate}
\section*{Chapitre~\ref{ch:contrat} --- Contrat et code}
\begin{enumerate}[label=\sffamily\bfseries\color{leaf}Q\arabic*.,leftmargin=2.2em]
  \item La référence \cmd{(SPEC-NNN)} de la spec dont il découle.
  \item Les entiers évitent les erreurs d'arrondi ; le FCFA n'a pas de centime en usage courant et se compte en entiers (\cmd{amountXof}).
  \item Poser la question ; ne pas deviner la règle.
\end{enumerate}
\section*{Chapitre~\ref{ch:recette} --- Recette}
\begin{enumerate}[label=\sffamily\bfseries\color{leaf}Q\arabic*.,leftmargin=2.2em]
  \item Specs \emph{Validée}, scénarios automatisés verts en CI, jeu de données anonymisé.
  \item Le responsable recette du client, jamais l'assistant.
  \item On la corrige sous 48\,h ouvrées : reproduire, chercher la cause racine, corriger le minimum, rejouer le scénario.
\end{enumerate}
\section*{Chapitre~\ref{ch:prod} --- Production}
\begin{enumerate}[label=\sffamily\bfseries\color{leaf}Q\arabic*.,leftmargin=2.2em]
  \item Le PV de recette signé par le client et aucune anomalie bloquante.
  \item Parce qu'un déclencheur vague ne se mesure pas : on ne sait pas quand l'appliquer.
  \item CI/CD, Docker, \cmd{.env}, migrations ou données de production, \cmd{git push}, service payant.
\end{enumerate}
\section*{Chapitre~\ref{ch:rapide} --- Cycle court}
\begin{enumerate}[label=\sffamily\bfseries\color{leaf}Q\arabic*.,leftmargin=2.2em]
  \item Par exemple : nouvelle dépendance, migration de données, plus d'une règle métier.
  \item Non : seul un humain la valide.
  \item Le travail a dépassé le cadre ; basculez vers le cycle complet.
\end{enumerate}
\section*{Chapitre~\ref{ch:git} --- Git}
\begin{enumerate}[label=\sffamily\bfseries\color{leaf}Q\arabic*.,leftmargin=2.2em]
  \item Les commits qui ne touchent que \cmd{docs/}, \cmd{.github/}, \cmd{scripts/}, \cmd{tests/} ou des fichiers \cmd{.md}, et ceux des branches \cmd{fix/}.
  \item Parce qu'il saute le hook local ; seule une CI côté serveur rend la règle absolue.
  \item En anglais.
\end{enumerate}
\section*{Chapitre~\ref{ch:apprendre} --- Apprentissage}
\begin{enumerate}[label=\sffamily\bfseries\color{leaf}Q\arabic*.,leftmargin=2.2em]
  \item Dans \cmd{apprentissage/}, pour ne pas toucher au vrai dossier de conception.
  \item La qualité du contenu : seulement la présence et le statut.
  \item Leçon~1 $\to$ \skill{forge-cadrage} ; 2 $\to$ \skill{forge-domaine} ; 3 $\to$ \skill{forge-spec} ; 4 $\to$ \skill{forge-architecture}.
\end{enumerate}

\end{document}
